\documentclass[11pt,reqno]{amsart}

% ================================================================
% Packages
% ================================================================
\usepackage{amsmath,amssymb,amsthm}
\usepackage{mathrsfs}
\usepackage{enumerate}
\usepackage[shortlabels]{enumitem}
\usepackage{booktabs}
\usepackage{array}
\usepackage{microtype}
\usepackage[dvipsnames]{xcolor}
\usepackage[colorlinks=true,linkcolor=blue!60!black,
 citecolor=green!40!black,urlcolor=blue!60!black]{hyperref}
\usepackage{tikz-cd}

% ================================================================
% Theorem environments
% ================================================================
\newtheorem{theorem}{Theorem}[section]
\newtheorem{proposition}[theorem]{Proposition}
\newtheorem{lemma}[theorem]{Lemma}
\newtheorem{corollary}[theorem]{Corollary}
\newtheorem{conjecture}[theorem]{Conjecture}
\theoremstyle{definition}
\newtheorem{definition}[theorem]{Definition}
\newtheorem{example}[theorem]{Example}
\newtheorem{construction}[theorem]{Construction}
\theoremstyle{remark}
\newtheorem{remark}[theorem]{Remark}
\newtheorem{warning}[theorem]{Warning}

% ================================================================
% Macros
% ================================================================
\newcommand{\cA}{\mathcal{A}}
\newcommand{\cB}{\mathcal{B}}
\newcommand{\cC}{\mathcal{C}}
\newcommand{\cD}{\mathcal{D}}
\newcommand{\cF}{\mathcal{F}}
\newcommand{\cH}{\mathcal{H}}
\newcommand{\cK}{\mathcal{K}}
\newcommand{\cL}{\mathcal{L}}
\newcommand{\cM}{\mathcal{M}}
\newcommand{\cO}{\mathcal{O}}
\newcommand{\cS}{\mathcal{S}}
\newcommand{\cV}{\mathcal{V}}
\newcommand{\cW}{\mathcal{W}}
\newcommand{\cZ}{\mathcal{Z}}
\newcommand{\bC}{\mathbb{C}}
\newcommand{\bN}{\mathbb{N}}
\newcommand{\bR}{\mathbb{R}}
\newcommand{\bZ}{\mathbb{Z}}
\newcommand{\barB}{\bar{B}}
\newcommand{\MC}{\mathrm{MC}}
\newcommand{\Sym}{\mathrm{Sym}}
\newcommand{\HS}{\mathrm{HS}}
\newcommand{\Vir}{\mathrm{Vir}}
\newcommand{\Res}{\mathrm{Res}}
\newcommand{\Tr}{\mathrm{Tr}}
\newcommand{\Hom}{\mathrm{Hom}}
\newcommand{\id}{\mathrm{id}}
\newcommand{\fg}{\mathfrak{g}}
\newcommand{\fH}{\mathfrak{H}}
\newcommand{\GG}{\Gamma}
\newcommand{\Dco}{D^{\mathrm{co}}}

\numberwithin{equation}{section}

% ================================================================
\begin{document}

\title[Analytic Sewing for Chiral Algebras]
{Analytic Sewing for Chiral Algebras}

\author{Raeez Lorgat}
\address{Perimeter Institute for Theoretical Physics,
 Waterloo, ON N2L 2Y5, Canada\newline
 Department of Physics \& Astronomy, University of Waterloo,
 Waterloo, ON N2L 3G1, Canada}
\email{rlorgat@perimeterinstitute.ca}

\date{\today}

\begin{abstract}
The bar complex of a chiral algebra is a purely algebraic
object: it encodes the operator product expansion through
residue-contracted point configurations on algebraic curves.
Partition functions, by contrast, require convergence of sewing
amplitudes across all genera.  This paper develops the analytic
completion that bridges the two.

We introduce the sewing envelope~$\cA^{\mathrm{sew}}$, the
universal locally convex completion for which every bordism
amplitude extends continuously, and formulate a single
Hilbert--Schmidt condition on the sector multiplication that
implies trace-class closed amplitudes at all genera.  For the
Heisenberg algebra we prove that the sewing envelope is
$\Sym A^2(D)$, the pair-of-pants amplitude is the second
quantisation of a one-particle Bergman restriction, and every
genus-$g$ partition function is a Fredholm determinant; this is
the decisive first case.  For lattice vertex algebras we prove
that the partition function factors into a Heisenberg Fredholm
determinant times a Siegel theta function, both absolutely
convergent on the Siegel upper half-space.  The general
HS-sewing theorem verifies the criterion for the entire standard
landscape: Heisenberg, affine Kac--Moody, Virasoro, $\cW_N$
algebras, and lattice vertex algebras.

We identify a three-layer analytic gap separating the algebraic
bar construction from convergent partition functions: the sewing
envelope for interacting algebras, metric independence of the
ind-Hilbert factorisation, and the coderived shadow at genus
$g \geq 1$.  The coderived category~$\Dco$ is the correct
receptacle: curved differentials ($d^2 = \kappa \cdot \omega_g
\neq 0$) are permitted, coacyclic complexes replace acyclic
ones, and the BV/bar comparison holds for all four shadow
classes including class~$\mathsf{M}$ (Virasoro, $\cW_N$),
where the naive chain-level identification fails.
\end{abstract}

\subjclass[2020]{17B69 (primary); 14H10, 30H20, 47B10, 81T40
 (secondary)}

\keywords{Chiral algebras, sewing, bar construction,
 Hilbert--Schmidt operators, Bergman space, Fredholm determinant,
 coderived category, partition functions}

\maketitle

\setcounter{tocdepth}{2}
\tableofcontents

% ================================================================
% 1. INTRODUCTION
% ================================================================

\section{Introduction}\label{sec:intro}

\subsection{The algebraic bar complex and the convergence
 problem}\label{subsec:problem}

The bar complex of a chiral algebra~$\cA$ is the tensor
coalgebra $B(\cA) = T^c(s^{-1}\bar\cA)$ with differential
extracted from the operator product expansion by residues at
collision diagonals in the Fulton--MacPherson
compactification~$\overline{\mathrm{FM}}_n(\bC)$.  The
construction is purely algebraic: it uses only the OPE
coefficients, the augmentation ideal~$\bar\cA = \ker(\varepsilon)$,
and the desuspension~$s^{-1}$ with~$|s^{-1}v| = |v| - 1$.
The bar differential satisfies $D^2 = 0$ on a fixed curve,
and the Maurer--Cartan element
$\Theta_\cA = D_\cA - d_0
\in \MC(\mathfrak{g}^{\mathrm{mod}}_\cA)$
governs the modular obstruction tower.

Partition functions require more.  A genus-$g$ partition
function is a trace
\begin{equation}\label{eq:partition-function}
Z_g(\cA;\,\Omega)
\;=\;
\Tr_{\cA}\bigl(q_1^{L_0 - c/24} \cdots q_g^{L_0 - c/24}\bigr)
\end{equation}
over a state space completed in a topology for which the
operator~$q^{L_0}$ is trace class.  The algebraic bar complex
produces formal power series in the sewing parameters; the
question is whether these series converge to honest complex
numbers.

This convergence problem is the analytic gap at the heart of the
modular Koszul duality programme~\cite{LorgatVolI}.  The five
theorems (A through D and~H) of that programme are stated and
proved algebraically; the analytic sewing package developed here
is the mechanism by which the algebraic skeleton extends to
convergent invariants.

\subsection{Main results}

We prove five results.

\paragraph{The sewing envelope.}
For any positive-energy chiral algebra~$\cA$, we construct
the \emph{sewing envelope}~$\cA^{\mathrm{sew}}$
(Definition~\ref{def:sewing-envelope}), the universal locally
convex completion for which every bordism amplitude extends
continuously, and establish its universal property
(Proposition~\ref{prop:sewing-universal}).

\paragraph{The Hilbert--Schmidt sewing condition.}
We formulate a single analytic hypothesis
(Definition~\ref{def:hs-sewing}): the sector multiplication
operators $m_{a,b}^c \colon H_a \otimes H_b \to H_c$ satisfy
\begin{equation}\label{eq:hs-intro}
\sum_{a,b,c \geq 0} q^{a+b+c}\,
\|m_{a,b}^c\|_{\HS}^2 \;<\; \infty
\end{equation}
for some $0 < q < 1$.  We prove
(Proposition~\ref{prop:closed-trace-class}) that this single
estimate delivers the full sewing package: every closed
amplitude is trace class on the weighted completion
$H_q := \widehat{\bigoplus}_n q^n H_n$.

\paragraph{The Heisenberg Fredholm determinant.}
For the Heisenberg algebra~$\cH_k$
(Theorem~\ref{thm:heisenberg-sewing}): the sewing envelope
is~$\Sym A^2(D)$; the pair-of-pants amplitude is the second
quantisation~$\GG(R)$ of a one-particle Bergman restriction;
every genus-$g$ partition function is the Fredholm determinant
\begin{equation}\label{eq:heisenberg-intro}
Z_g(\cH_k;\,\Omega)
\;=\;
(\det\operatorname{Im}\Omega)^{-k/2}
\cdot (\det{}'_\zeta \Delta_{\Sigma_g})^{-k/2}.
\end{equation}

\paragraph{Lattice sewing.}
For a lattice vertex algebra~$V_\Lambda$ of rank~$r$
(Theorem~\ref{thm:lattice-sewing}): the genus-$g$ partition
function factors as
\begin{equation}\label{eq:lattice-intro}
Z_g(V_\Lambda;\,\Omega)
\;=\;
Z_g(\cH_r;\,\Omega)
\cdot \Theta_\Lambda(\Omega),
\end{equation}
where~$\Theta_\Lambda(\Omega)
= \sum_{\lambda \in \Lambda^g}
e^{i\pi\lambda^T\Omega\lambda}$ is the Siegel theta function
of the lattice, absolutely convergent on the Siegel upper
half-space~$\fH_g$.

\paragraph{The general HS-sewing theorem.}
Any positive-energy chiral algebra with subexponential sector
growth and polynomial OPE coefficient growth satisfies the
HS-sewing condition
(Theorem~\ref{thm:general-hs-sewing}).  The criterion verifies
for the Heisenberg, affine Kac--Moody at non-critical level,
Virasoro at every central charge, $\cW_N$ algebras, and lattice
vertex algebras.

\subsection{The three-layer gap and the coderived
 category}\label{subsec:gap-intro}

Convergent partition functions are the beginning of the analytic
theory, not its end.  Three layers of structure separate the
algebraic bar construction from a complete analytic modular
Koszul programme:
\begin{enumerate}[label=\textup{(\arabic*)}]
\item \emph{The sewing envelope for interacting algebras.}
 For the Heisenberg and lattice algebras, the sewing envelope
 admits an explicit description (Bergman symmetric algebra,
 charge-refined Hilbert completion).  For non-abelian chiral
 algebras (affine Kac--Moody at generic level, Virasoro,
 $\cW_N$), the sewing envelope exists abstractly but lacks a
 concrete Hilbert-space identification.
\item \emph{Metric independence of the ind-Hilbert
 factorisation.}  The ind-Hilbert completion of
 Moriwaki~\cite{Moriwaki26b} depends a~priori on a conformal
 metric on the curve; the algebraic bar complex does not.
 Identifying the two requires proving that the analytic
 completion is independent of the metric, a functional-analytic
 problem with no algebraic counterpart.
\item \emph{The coderived shadow at $g \geq 1$.}  At genus
 $g \geq 1$ the bar complex carries curvature:
 $d_{\mathrm{fib}}^2 = \kappa(\cA) \cdot \omega_g \neq 0$.
 The ordinary derived category is empty (every object is
 acyclic).  The correct receptacle is the coderived
 category~$\Dco$, defined by totalising the curvature bicomplex.
\end{enumerate}

In Sections~\ref{sec:gap} and~\ref{sec:coderived} we develop
the third layer in detail: the coderived category as the natural
home for curved bar complexes, the notion of coacyclicity
replacing acyclicity, and the BV/bar comparison theorem in~$\Dco$
for all four shadow classes (G, L, C, M) of the standard landscape.

\subsection{Relation to prior work}\label{subsec:prior}

Segal~\cite{Segal88} formulated the sewing axiom as the
defining property of a conformal field theory.
Zhu~\cite{Zhu96} proved modular invariance of genus-1 characters
for rational vertex operator algebras.
Huang~\cite{Huang05} extended the sewing of sphere amplitudes
to genus one and by induction to all genera, under a list of
analytic convergence hypotheses.  The HS criterion of this paper
is strictly stronger: a single norm
estimate~\eqref{eq:hs-intro} implies all convergences
simultaneously.

Moriwaki~\cite{Moriwaki26b} identified the Heisenberg sewing
envelope with~$\Sym A^2(D)$ through factorisation homology
valued in ind-Hilbert spaces.  We reinterpret this as a
one-particle reduction of the pair-of-pants operator and deduce
the Fredholm determinant
formula~\eqref{eq:heisenberg-intro}.

The Belavin--Knizhnik theorem~\cite{BK86} expresses the bosonic
string measure as a product of determinants of Laplacians; our
formula~\eqref{eq:heisenberg-intro} is the vertex-algebraic
incarnation.  Beilinson--Drinfeld~\cite{BD04} develop the
algebraic theory of chiral algebras and factorisation; our bar
complex is the explicit chain model of their factorisation
coalgebra.

The coderived category for curved dg-modules was introduced
by Positselski~\cite{Positselski11}; its application to
curved bar complexes in the chiral setting appears
in~\cite{LorgatVolI}.  The BV/bar identification in the
coderived category (Theorem~\ref{thm:bv-bar-coderived})
extends the genus-0 chain-level result of Costello and
Gwilliam~\cite{CG17} to all genera.

\subsection{The modular Koszul duality
 context}\label{subsec:programme-context}

The analytic sewing programme is one component of a larger
algebraic structure.  The modular Koszul duality programme
of~\cite{LorgatVolI} constructs, for every chirally Koszul
algebra~$\cA$, a package of algebraic invariants:
\begin{enumerate}[label=\textup{(\alph*)}]
\item the bar complex
 $B(\cA) = T^c(s^{-1}\bar\cA)$ with differential and
 deconcatenation coproduct;
\item the Maurer--Cartan element
 $\Theta_\cA = D_\cA - d_0
 \in \MC(\mathfrak{g}^{\mathrm{mod}}_\cA)$;
\item the shadow obstruction tower
 $\Theta_\cA = \sum_{r \geq 2} S_r(\cA)$
 with $S_2 = \kappa(\cA)$ the modular characteristic;
\item the Koszul dual~$\cA^!$ and the complementarity
 relation $\kappa(\cA) + \kappa(\cA^!) = K(\cA)$;
\item five theorems (A through D and~H) controlling the
 structure of the bar complex, bar-cobar inversion,
 complementarity, the genus expansion, and chiral Hochschild
 cohomology.
\end{enumerate}

All of this is purely algebraic: the bar complex is built from
OPE coefficients and residues; no convergence hypothesis enters.
The analytic sewing programme begins where the algebraic
programme ends: it asks when the formal genus expansion
$\Theta_\cA = \sum_g \Theta_\cA^{(g)}$ converges to a
well-defined partition function~$Z_g(\cA;\,\Omega)$.

The algebraic programme establishes the formula; the analytic
programme establishes the number.  The two are connected by the
HS-sewing condition, which is precisely the analytic hypothesis
that converts the algebraic genus expansion into a convergent
Fredholm expansion.

\begin{remark}[The four shadow classes and
 sewing]\label{rem:shadow-sewing-intro}
The shadow classification (G, L, C, M) of the standard
landscape controls the complexity of the sewing programme:
\begin{enumerate}[label=\textup{(\alph*)}]
\item \emph{Class~G} (Heisenberg, lattice): the shadow tower
 terminates at degree~$2$; the sewing programme is
 fully explicit (Sections~\ref{sec:heisenberg}
 and~\ref{sec:lattice}).
\item \emph{Class~L} (affine Kac--Moody): the shadow tower
 terminates at degree~$3$; the sewing programme is proved
 abstractly (HS-sewing, Theorem~\ref{thm:general-hs-sewing})
 but the sewing envelope is not yet identified concretely.
\item \emph{Class~C} ($\beta\gamma$, $bc$): the shadow tower
 terminates at degree~$4$; the BV/bar comparison is
 conditional on harmonic decoupling.
\item \emph{Class~M} (Virasoro, $\cW_N$): the shadow tower is
 infinite; the chain-level BV/bar comparison fails; the
 coderived category is essential
 (Section~\ref{sec:coderived}).
\end{enumerate}
\end{remark}

\subsection{Outline}

Section~\ref{sec:sewing-envelope} constructs the sewing
envelope.  Section~\ref{sec:hs-criterion} formulates the
Hilbert--Schmidt sewing condition and proves the trace-class
theorem.  Section~\ref{sec:heisenberg} treats the decisive
first case: the Heisenberg Fredholm determinant.
Section~\ref{sec:lattice} develops lattice sewing via theta
functions.  Section~\ref{sec:general-criterion} proves the
general HS-sewing theorem and verifies the criterion on the
standard landscape.  Section~\ref{sec:gap} analyses the
three-layer analytic gap.  Section~\ref{sec:coderived}
develops the coderived category as the correct receptacle for
curved bar complexes at $g \geq 1$.

\subsection{Conventions}\label{subsec:conventions}

All chiral algebras are $\bZ_{\geq 0}$-graded
(positive energy), with finite-dimensional weight spaces.
Grading is cohomological: $|d| = +1$.  Desuspension lowers
degree: $|s^{-1}v| = |v| - 1$.  The bar complex uses the
augmentation ideal: $B(\cA) = T^c(s^{-1}\bar\cA)$ with
$\bar\cA = \ker(\varepsilon)$.  The Dedekind eta function is
$\eta(\tau) = q^{1/24}\prod_{n \geq 1}(1-q^n)$ with
$q = e^{2\pi i \tau}$.


% ================================================================
% 2. THE SEWING ENVELOPE
% ================================================================

\section{The sewing envelope}\label{sec:sewing-envelope}

\subsection{Algebraic cores and bordism
 amplitudes}\label{subsec:algebraic-cores}

An \emph{algebraic chiral core} is a
$\bZ_{\geq 0}$-graded vector space
$\cA_{\mathrm{alg}} = \bigoplus_{n \geq 0} H_n$
equipped with the operator product expansion of a vertex
operator algebra of positive energy, with
$H_0 = \bC \cdot \mathbf{1}$ and $\dim H_n < \infty$.
The sector components~$H_n$ are the conformal-weight
eigenspaces.  The augmentation
$\varepsilon \colon \cA_{\mathrm{alg}} \to \bC$ is the
projection to~$H_0$.  The augmentation ideal
$\bar\cA = \ker(\varepsilon)$ is the input to the bar
construction $B(\cA) = T^c(s^{-1}\bar\cA)$.

A finite conformally flat bordism is a compact Riemann
surface~$\Sigma$ with parametrised in/out boundary collars.
Each such~$\Sigma$ determines an algebraic amplitude
\[
A_\Sigma^{\mathrm{alg}} \colon
\cA_{\mathrm{alg}}^{\otimes m}
\longrightarrow
\cA_{\mathrm{alg}}^{\otimes \ell},
\]
a~priori valued in formal series of conformal-weight
components.  Matrix coefficients
\[
\langle \eta,\,
A_\Sigma^{\mathrm{alg}}
(\xi_1 \otimes \cdots \otimes \xi_m)
\rangle
\;\in\;
\bC[[\text{moduli of } \Sigma]]
\]
are the objects whose convergence the sewing programme must
establish.

\subsection{Construction of the sewing
 envelope}\label{subsec:sewing-envelope-construction}

\begin{definition}[Sewing envelope]\label{def:sewing-envelope}
For every finite conformally flat bordism~$\Sigma$ and every
choice of boundary vectors
$\xi_i \in \cA_{\mathrm{alg}}$,
$\eta \in \cA_{\mathrm{alg}}^\vee$, define the seminorm
\begin{equation}\label{eq:sewing-seminorm}
p_{\Sigma,\xi,\eta}(a)
\;=\;
\bigl|\langle \eta,\,
A_\Sigma^{\mathrm{alg}}
(\xi_1 \otimes \cdots \otimes a \otimes \cdots \otimes \xi_m)
\rangle\bigr|.
\end{equation}
The \emph{sewing envelope}~$\cA^{\mathrm{sew}}$ is the
Hausdorff completion of~$\cA_{\mathrm{alg}}$ for the locally
convex topology generated by all
seminorms~$p_{\Sigma,\xi,\eta}$.
\end{definition}

The sewing envelope is the smallest completion that remembers
every local-to-global operation at once: it is tautologically
the correct target for any analytic bootstrap.

\begin{proposition}[Universal property]
\label{prop:sewing-universal}
Let $E$ be any complete locally convex realisation of
$\cA_{\mathrm{alg}}$ for which every algebraic amplitude
$A_\Sigma^{\mathrm{alg}}$ extends continuously.  Then the
identity on~$\cA_{\mathrm{alg}}$ extends uniquely to a
continuous map $\cA^{\mathrm{sew}} \to E$.
\end{proposition}

\begin{proof}
By hypothesis, each seminorm~$p_{\Sigma,\xi,\eta}$ is
continuous on~$E$.  Hence the identity
$\cA_{\mathrm{alg}} \to E$ is continuous for the sewing
topology.  By completeness of~$E$, it extends uniquely from
the Hausdorff completion~$\cA^{\mathrm{sew}}$.
\end{proof}

\begin{remark}[Comparison with other completions]
\label{rem:completion-comparison}
The sewing envelope need not coincide with naive completions:
the $q$-adic completion
$\widehat{\bigoplus}_n q^n H_n$, the Hilbert-space completion
with respect to a fixed invariant form, Zhu's $C_2$-quotient,
or Moriwaki's ind-Hilbert completion.  All of these are quotients
or subspaces of~$\cA^{\mathrm{sew}}$ under the universal
property, but none is a~priori large enough to carry every
bordism amplitude as a continuous map.  The HS-sewing criterion
of the next section is the device by which one identifies a
concrete Hilbert completion
$H_q \subset \cA^{\mathrm{sew}}$.
\end{remark}

\subsection{Functoriality}\label{subsec:sewing-functoriality}

The sewing envelope construction is functorial in the algebraic
core.  A morphism of vertex algebras
$f \colon \cA \to \cB$ that preserves the conformal grading
and the OPE induces a continuous map
$f^{\mathrm{sew}} \colon \cA^{\mathrm{sew}} \to
\cB^{\mathrm{sew}}$.  In particular, any subalgebra
inclusion~$\cA \hookrightarrow \cB$ (such as the inclusion of
a Heisenberg subalgebra into a lattice vertex algebra) induces
a continuous inclusion of sewing envelopes.

\begin{proposition}[Direct-sum decomposition]
\label{prop:sewing-direct-sum}
Let $\cA = \bigoplus_{\gamma \in I} \cA_\gamma$ be a
decomposition of the algebraic core into sectors preserved by
the pair-of-pants multiplication.  Then
\[
\cA^{\mathrm{sew}}
\;=\;
\widehat{\bigoplus}_{\gamma \in I}\,
\cA_\gamma^{\mathrm{sew}},
\]
where the completion is in the product topology.
\end{proposition}

\begin{proof}
Each seminorm~$p_{\Sigma,\xi,\eta}$ decomposes by sector
selection (insertion of projection operators onto the
sector~$\gamma$ at each boundary circle); the cross-terms are
separate seminorms.  The completion of a direct sum in a product
of locally convex topologies is the product of completions.
\end{proof}


\subsection{The five levels of analytic
 structure}\label{subsec:five-levels}

The sewing envelope sits at the fourth of five levels of
analytic structure for chiral algebras, arranged in a strict
chain of implications:

\begin{theorem}[Five-level hierarchy]
\label{thm:five-levels}
The following conditions on a positive-energy chiral
algebra~$\cA$ form a strict chain:
\begin{enumerate}[label=\textup{(\roman*)}]
\item \textup{(\textbf{Algebraic})}\;
 $D_\cA^2 = 0$ on the bar complex.  Universal: holds for
 every chiral algebra.
\item \textup{(\textbf{Coderived})}\;
 The homotopy transfer theorem applies weight-by-weight in
 the coderived category~$\Dco$.  Requires: positive energy
 ($\dim H_n < \infty$).
\item \textup{(\textbf{Koszul})}\;
 The curved bar complex and the flat model are linked by the
 coderived-to-derived comparison.  Requires: PBW collapse at
 the $E_2$ page of the bar spectral sequence.
\item \textup{(\textbf{Analytic})}\;
 The genus-$g$ partition functions converge absolutely.
 Requires: HS-sewing (polynomial OPE growth, subexponential
 sector growth).
\item \textup{(\textbf{Modular functor})}\;
 Conformal blocks are finite-dimensional and the modular group
 acts.  Requires: $C_2$-cofiniteness.
\end{enumerate}
The first four levels form the chain
$\textup{(iv)} \subsetneq \textup{(iii)} \subsetneq
\textup{(ii)} \subsetneq \textup{(i)}$.
Level~$\textup{(v)}$ implies
$\textup{(iv)}$, $\textup{(ii)}$, and $\textup{(i)}$, but
its relation to~$\textup{(iii)}$ is not one of containment:
$C_2$-cofinite simple quotients~$L_k(\fg)$ at admissible level
satisfy~$\textup{(v)}$ but may fail~$\textup{(iii)}$, while
Koszul universal algebras~$V^k(\fg)$ at generic level
satisfy~$\textup{(iii)}$ but not~$\textup{(v)}$.
\end{theorem}

\begin{proof}[Proof sketch]
Level~(i): $D^2 = 0$ holds because the bar differential is the
Feynman-transform differential of a modular operad algebra, and
$\partial^2 = 0$ on~$\overline\cM_{g,n}$ implies $D^2 = 0$.
Level~(ii): the positive-energy axiom gives weight-by-weight
finite-dimensional SDR data for the homotopy transfer.
Level~(iii): the PBW collapse gives diagonal Ext vanishing.
Level~(iv): HS-sewing (Theorem~\ref{thm:general-hs-sewing}).
Level~(v): Zhu's theorem~\cite{Zhu96} gives finite-dimensional
conformal blocks and modular group action.

The inclusions are strict.  For (i)$\neq$(ii): a chiral algebra
without positive energy satisfies $D^2 = 0$ but lacks the SDR.
For (ii)$\neq$(iii): admissible-level quotients where bar-Ext
differs from ordinary-Ext.
For (iii)$\neq$(iv): a chirally Koszul algebra with
super-polynomial OPE growth.
For (iv)$\neq$(v): the universal $\cW$-algebra $V^k(\fg)$ at
generic irrational level.
\end{proof}

\begin{remark}[The standard landscape]
\label{rem:standard-landscape-levels}
The standard landscape (Heisenberg, affine Kac--Moody at
non-critical level, Virasoro, $\cW_N$, $\beta\gamma$ and $bc$
systems, lattice vertex algebras) satisfies levels
(i) through~(iv).  Those families that are also
$C_2$-cofinite (simple quotients at admissible levels,
self-dual lattice VOAs) satisfy all five levels.  The sewing
envelope is the universal object at level~(iv); the present
paper develops this level in detail.
\end{remark}

\subsection{The bar complex as a family over
 moduli}\label{subsec:bar-over-moduli}

The bar complex at genus~$g$ is a family of cochain complexes
parametrised by~$\overline\cM_g$.  On the smooth
locus~$\cM_g$ the differential satisfies $d^2 = 0$; on the
compactification~$\overline\cM_g$ the differential acquires
curvature $d_{\mathrm{fib}}^2 = \kappa \cdot \omega_g$ from
the boundary strata.

The Gauss--Manin connection
$\nabla^{\mathrm{GM}} \colon B^{(g)}(\cA) \to
B^{(g)}(\cA) \otimes \Omega^1_{\overline\cM_g}$
controls the variation of the bar complex with moduli.  The
total differential
$D_g = d_{\mathrm{fib}} + \nabla^{\mathrm{GM}}$ satisfies
$D_g^2 = 0$: the fiberwise curvature
$d_{\mathrm{fib}}^2 = \kappa\omega_g$ is exactly cancelled
by the Gauss--Manin curvature
$(\nabla^{\mathrm{GM}})^2 = -\kappa\omega_g$.

The analytic sewing programme extends this family to a family of
Hilbert spaces: on each fiber (a fixed Riemann surface), the
algebraic bar complex is completed to a Hilbert space by the
HS-sewing condition, and the Gauss--Manin connection extends
continuously to the completed fibers.

\begin{proposition}[Continuity of the Gauss--Manin connection]
\label{prop:gm-continuity}
Under HS-sewing, the Gauss--Manin connection
$\nabla^{\mathrm{GM}}_g$ extends to a bounded operator on the
weighted completion~$H_q$ for each~$g$, and the total
differential~$D_g = d_{\mathrm{fib}} + \nabla^{\mathrm{GM}}_g$
satisfies $D_g^2 = 0$ on the completed bar complex.
\end{proposition}

\begin{proof}
The Gauss--Manin connection acts on bar elements by variation
of the bar propagator $d\log E(z,w)$ with respect to moduli.
The variation is a holomorphic $(1,0)$-form on~$\overline\cM_g$
valued in operators on~$H_q$; the HS bound on the pair-of-pants
multiplication controls the operator norm of these
variations, because each variation is expressed as a residue of
the bar differential applied to a collar region.  The
identity~$D_g^2 = 0$ is algebraic (it holds on the dense
algebraic core) and extends by continuity.
\end{proof}


% ================================================================
% 3. THE HILBERT--SCHMIDT SEWING CONDITION
% ================================================================

\section{The Hilbert--Schmidt sewing condition}
\label{sec:hs-criterion}

The universal sewing envelope of
Definition~\ref{def:sewing-envelope} is an abstract object.
To compute with it one needs a concrete Hilbert completion on
which the sewing operators are not merely defined but small
enough in norm to compose.

\subsection{The condition}\label{subsec:hs-definition}

\begin{definition}[HS-sewing]\label{def:hs-sewing}
Let $\cA_{\mathrm{alg}} = \bigoplus_n H_n$ be an algebraic
chiral core with a pre-Hilbert structure on each sector
(for example, the restriction of the invariant bilinear form
to~$H_n$).  Write the pair-of-pants sector multiplications as
\[
m_{a,b}^c \colon H_a \otimes H_b \longrightarrow H_c.
\]
Say $\cA$ satisfies \emph{HS-sewing} at parameter
$q \in (0,1)$ if
\begin{equation}\label{eq:hs-sewing}
\sum_{a,b,c \geq 0} q^{a+b+c}\,
\|m_{a,b}^c\|_{\HS}^2
\;<\; \infty.
\end{equation}
The \emph{weighted completion} is
$H_q := \widehat{\bigoplus}_{n \geq 0} q^n H_n$, with inner
product
$\langle v, w \rangle_q = q^{2n}\langle v, w\rangle_{H_n}$
for $v, w \in H_n$.
\end{definition}

The scaling~$q^{2n}$ is the unique convention under which the
weighted pair-of-pants~$m_q$ is Hilbert--Schmidt precisely
when~\eqref{eq:hs-sewing} holds.

\subsection{The trace-class theorem}
\label{subsec:trace-class}

\begin{proposition}[Closed amplitudes are trace class]
\label{prop:closed-trace-class}
Suppose $\cA$ satisfies HS-sewing at parameter~$q$.  Then:
\begin{enumerate}[label=\textup{(\roman*)}]
\item the weighted pair-of-pants
$m_q \colon H_q \otimes H_q \to H_q$ is Hilbert--Schmidt, with
\[
\|m_q\|_{\HS}^2
\;=\;
\sum_{a,b,c} q^{a+b+c}\, \|m_{a,b}^c\|_{\HS}^2;
\]
\item every closed amplitude of~$\cA$, built from any finite
composition of pair-of-pants multiplications and traces, is
trace class on~$H_q$;
\item the sewing envelope~$\cA^{\mathrm{sew}}$ admits a
continuous linear map
$\cA^{\mathrm{sew}} \to H_q$ with dense image.
\end{enumerate}
\end{proposition}

\begin{proof}
Claim~(i): identifying $H_a \otimes H_b$ with the
weight-$(a+b)$ subspace of $H_q^{\widehat{\otimes} 2}$,
the sector component $m_{a,b}^c$ contributes
$q^a \cdot q^b \cdot q^{-c} \cdot q^{2c} = q^{a+b+c}$
to the Hilbert--Schmidt norm of~$m_q$ (the factor~$q^{-c}$
comes from the output side, the factor~$q^{2c}$ from the
normalisation of the inner product on~$H_c$).  Summing gives
$\|m_q\|_{\HS}^2 =
\sum q^{a+b+c}\|m_{a,b}^c\|_{\HS}^2 < \infty$.

Claim~(ii): a genus-$g$ Riemann surface with~$n$ punctures
admits a pants decomposition into $2g - 2 + n$ pairs of pants.
Each internal sewing circle corresponds to a Hilbert--Schmidt
pair-of-pants composition, and the composition of two
Hilbert--Schmidt operators is trace class.  Every closed
amplitude (no free external legs) evaluates to the trace of
a finite composition of trace-class operators.

Claim~(iii): the seminorms~$p_{\Sigma,\xi,\eta}$ of
Definition~\ref{def:sewing-envelope} are all dominated by
the~$H_q$ inner product, because each matrix coefficient
factorises through the trace-class operator of claim~(ii).
Continuity of the identity
$\cA_{\mathrm{alg}} \to H_q$ for the sewing topology follows.
\end{proof}

\begin{corollary}[Partition functions converge]
\label{cor:partition-convergence}
Under HS-sewing at parameter~$q$, the genus-$g$ partition
function
$Z_g(\cA;\,\Omega)
= \Tr(q_1^{L_0 - c/24}\cdots q_g^{L_0 - c/24})$
converges absolutely for all period matrices
$\Omega \in \fH_g$ with
$|q_j| = |e^{2\pi i \Omega_{jj}}| < q$, producing a
holomorphic function on a domain in the Siegel upper
half-space.
\end{corollary}

\begin{proof}
The partition function is a closed amplitude (trace over the
state space with $g$ sewing circles).  By
Proposition~\ref{prop:closed-trace-class}(ii), it is the trace
of a trace-class operator on~$H_q$, hence absolutely convergent.
Holomorphicity in the moduli follows from the absolute convergence
and the holomorphic dependence of the sewing kernel on~$\Omega$.
\end{proof}

\begin{corollary}[Gram positivity]
\label{cor:hs-positivity}
Under HS-sewing, the connected sewing amplitudes
$a_\cA(N) \in \bR_{\geq 0}$ define a positive-semidefinite
inner product on generating Dirichlet polynomials.
\end{corollary}

\begin{proof}
Trace-class operators form a two-sided ideal with a well-defined
trace, positive when applied to a positive operator; the sewing
amplitude is a composition of positive pair-of-pants operators,
and positivity passes through.
\end{proof}

\subsection{The sewing Gram matrix and its
 determinant}\label{subsec:gram-matrix}

The Hilbert--Schmidt bound controls not only convergence but
spectral asymptotics.  Let $K_g$ denote the sewing operator
(composition of pair-of-pants restrictions around a cycle of
the genus-$g$ surface).  The connected free energy is
\begin{equation}\label{eq:connected-free-energy}
F_g^{\mathrm{conn}}(\cA;\,q)
\;=\;
-\log \det(1 - K_g),
\end{equation}
convergent whenever $\|K_g\|_1 < 1$.  The eigenvalues
$\lambda_1 \geq \lambda_2 \geq \cdots$ of~$K_g$ satisfy
$\sum_j \lambda_j < \infty$ (trace class), and the Fredholm
determinant $\det(1 - K_g) = \prod_j(1 - \lambda_j)$ converges
absolutely.  For the Heisenberg, the eigenvalues are $q^n$
($n \geq 1$) and the Fredholm determinant is the Dedekind
eta function (see Section~\ref{sec:heisenberg}).


\subsection{The HS-sewing condition and the shadow
 classes}\label{subsec:hs-shadow}

The HS-sewing condition interacts with the shadow classification
(G, L, C, M) of the modular Koszul duality programme.  The
shadow class determines the OPE pole structure and hence the
polynomial degree~$N$ in the OPE growth bound.

\begin{proposition}[Shadow class and OPE growth]
\label{prop:shadow-hs}
The OPE growth exponent~$N$ in the bound
$|C^{c,k}_{a,i;\,b,j}| \leq K(a+b+c+1)^N$ is controlled
by the shadow class:
\begin{center}
\renewcommand{\arraystretch}{1.15}
\begin{tabular}{llcl}
\toprule
Class & Families & $N$ & Mechanism \\
\midrule
$\mathsf{G}$ & Heisenberg, free fields
 & $1$ & quadratic OPE \\
$\mathsf{L}$ & affine Kac--Moody
 & $1$ & linear current algebra \\
$\mathsf{C}$ & $\beta\gamma$, $bc$
 & $1$ & first-order systems \\
$\mathsf{M}$ & Virasoro, $\cW_N$
 & $2N$ & higher-spin fields \\
\bottomrule
\end{tabular}
\end{center}
Classes~G, L, and~C all have $N = 1$ (linear OPE growth);
class~M has $N$ growing with the number of generators.
\end{proposition}

\begin{proof}
For class~G: the OPE $J(z)J(w) \sim k/(z-w)^2$ has no field
on the right-hand side beyond the central term; the structure
constants are bounded by the level~$k$, independent of weight.

For class~L: the OPE $J^a(z)J^b(w) \sim
k\delta^{ab}/(z-w)^2 + f^{ab}{}_cJ^c(w)/(z-w)$
has at most one derivative insertion, giving $|C| \leq K(n+1)$.

For class~C: the $\beta\gamma$ system has OPE
$\beta(z)\gamma(w) \sim 1/(z-w)$, a simple pole; all
higher correlators factor through this pairing by Wick's
theorem.

For class~M: the Virasoro OPE is quartic, and normal-ordered
products of~$T$ with itself produce structure constants growing
as $(n+1)^2$.  For $\cW_N$, the spin-$s$ generator
contributes $2s$ poles, and the OPE growth is bounded by
$(n+1)^{2N}$.
\end{proof}

The shadow class also determines the convergence rate of the
HS-sewing sum.  For classes~G, L, C (all with $N = 1$),
the HS-sewing sum converges rapidly:
\[
\sum_{a,b,c} q^{a+b+c}\|m_{a,b}^c\|_{\HS}^2
\;\leq\;
9K^2\Bigl(\sum_n q^n \dim H_n (n+1)^2\Bigr)
\Bigl(\sum_n q^n \dim H_n\Bigr)^2,
\]
and the rate-limiting factor is the sector growth
$\dim H_n$.  For class~M with large~$N$, the polynomial
factor $(n+1)^{2N}$ slows convergence but does not affect
the qualitative conclusion (the exponential decay of~$q^n$
dominates any polynomial growth).

\subsection{Quantitative HS bounds}\label{subsec:quantitative-hs}

The HS-sewing condition provides not just qualitative
convergence but quantitative bounds on partition functions.

\begin{proposition}[Partition function bound]
\label{prop:partition-bound}
Let $\cA$ satisfy HS-sewing at parameter~$q$.  Then the
genus-$g$ partition function satisfies
\begin{equation}\label{eq:partition-bound}
|Z_g(\cA;\,\Omega)|
\;\leq\;
\bigl(\|m_q\|_{\HS}\bigr)^{2(2g-2)}
\end{equation}
for all $\Omega \in \fH_g$ with
$|q_j| < q$.
\end{proposition}

\begin{proof}
A pants decomposition of the genus-$g$ surface has
$2g - 2$ pairs of pants and $3g - 3$ sewing circles.
Each pair of pants contributes a factor of
$\|m_q\|_{\HS}$ (the Hilbert--Schmidt norm of the
pair-of-pants multiplication); each sewing circle
composes two such factors.  The closed amplitude (the
trace) is bounded by the product of all $\HS$ norms:
$|Z_g| \leq \|m_q\|_{\HS}^{2(2g-2)}$ (the exponent
counts the number of internal pair-of-pants composition
operations in a genus-$g$ surface with no punctures).
\end{proof}

\begin{corollary}[Growth in genus]\label{cor:growth-genus}
Under HS-sewing, the partition functions grow at most
exponentially in genus:
$\log|Z_g| \leq (2g-2)\log\|m_q\|_{\HS}^2$.
For the Heisenberg,
$\|m_q\|_{\HS}^2 = k\sum_n q^{2n+1}/(1-q)^2$, and the
growth rate is controlled by the level~$k$ and the sewing
parameter~$q$.
\end{corollary}

This exponential-in-genus bound is consistent with the
predictions of the Belavin--Knizhnik factorisation~\cite{BK86}
and with the genus-$g$ Faber--Pandharipande
evaluation~$\lambda_g^{\mathrm{FP}}$, which grows as
$(2g)!/(2\pi)^{2g}$.

\subsection{Optimality of the HS condition}
\label{subsec:hs-optimality}

The HS-sewing condition is not sharp: it is a sufficient
condition for convergent partition functions, not a necessary
one.  Two weaker conditions that also imply convergence in
special cases:

\begin{enumerate}[label=\textup{(\alph*)}]
\item \emph{$C_2$-cofiniteness:} Zhu's condition~\cite{Zhu96}
 implies finite-dimensional conformal blocks and hence
 convergent characters.  This is strictly stronger than
 HS-sewing for rational VOAs (it implies modularity, which
 HS-sewing does not), but strictly weaker in scope (it applies
 only to $C_2$-cofinite algebras, while HS-sewing applies to
 generic-level universal algebras).
\item \emph{Trace-class pair-of-pants:} one could replace the
 Hilbert--Schmidt condition on~$m_{a,b}^c$ with a trace-class
 condition on the composed operator
 $m_q \circ m_q^*$.  This is strictly weaker than HS-sewing
 (trace class $\subset$ Hilbert--Schmidt) but harder to verify
 from OPE data alone.
\end{enumerate}

The HS condition strikes the optimal balance between
generality (it applies to the full standard landscape) and
verifiability (it can be checked from character formulas and
OPE tables).


% ================================================================
% 4. THE HEISENBERG SEWING: FREDHOLM DETERMINANT
% ================================================================

\section{The Heisenberg sewing: Fredholm determinant}
\label{sec:heisenberg}

The Heisenberg vertex algebra~$\cH_k$ is the decisive first
case: a chiral algebra for which the entire analytic sewing
programme admits a complete, explicit description.  Every
step reduces to Fock-space combinatorics and Bergman-space
functional analysis.

\subsection{The Heisenberg algebra}\label{subsec:heisenberg-ope}

The Heisenberg vertex algebra~$\cH_k$ at level~$k$ is
generated by a single current~$J(z)$ with operator product
expansion
\[
J(z)\, J(w)
\;\sim\;
\frac{k}{(z-w)^2}.
\]
The sectors~$H_n$ are spanned by monomials
$a_{-n_1-1}\cdots a_{-n_j-1}\mathbf{1}$ with
$\sum(n_i + 1) = n$; the sector generating function is
$\sum_{n \geq 0} \dim H_n\,q^n
= \prod_{n \geq 1}(1 - q^n)^{-1}$,
equivalently $q^{1/24}/\eta(q)$.  The Heisenberg is
the simplest member of the standard landscape, class~G in the
shadow classification (the shadow tower terminates at
degree~$2$), with $\kappa(\cH_k) = k$.

\subsection{The Heisenberg OPE and HS-sewing
 verification}\label{subsec:heisenberg-hs-verification}

The HS-sewing condition for the Heisenberg is verified
directly from the OPE.  The pair-of-pants multiplication
$m_{a,b}^c \colon H_a \otimes H_b \to H_c$ is nonzero
only when $c \leq a + b$ (the weight cannot increase beyond
the total input weight).  The OPE
$J(z)J(w) \sim k/(z-w)^2$ produces structure constants
bounded by~$k$: the $n$-th mode $a_n$ satisfies
$[a_m, a_n] = km\delta_{m+n,0}$, so the only nontrivial
sector multiplication is the Wick contraction
$m_{a,b}^{a+b-1} \colon
a_{-a-1}\mathbf{1} \otimes a_{-b-1}\mathbf{1}
\mapsto k(a+1)\delta_{b,0}\,a_{-a-1}\mathbf{1}$
(plus normal-ordered products).  The Hilbert--Schmidt norm
is
\[
\|m_{a,b}^c\|_{\HS}^2
\;\leq\;
k^2 p(a)p(b)p(c)(a+b+c+1)^2,
\]
and the HS-sewing sum is
\[
\sum_{a,b,c} q^{a+b+c}\|m_{a,b}^c\|_{\HS}^2
\;\leq\;
9k^2\Bigl(\sum_n q^n p(n)(n+1)^2\Bigr)
\Bigl(\sum_n q^n p(n)\Bigr)^2.
\]
Each series converges for $0 < q < 1$ because $p(n)$ grows
subexponentially ($\log p(n) \sim \pi\sqrt{2n/3}$) while
$q^n$ decays exponentially.  The radius of convergence
$q = 1$ is sharp: at $q = 1$ the series $\sum p(n) = \infty$.

\begin{remark}[The partition number and subexponential growth]
\label{rem:partition-growth}
The Hardy--Ramanujan asymptotic
$p(n) \sim \frac{1}{4n\sqrt{3}}
\exp\bigl(\pi\sqrt{2n/3}\bigr)$
shows that $\log p(n) = O(\sqrt{n})$, hence
$p(n) \leq C\exp(\alpha\sqrt{n})$ for appropriate~$C$
and~$\alpha$.  The exponential decay
$q^n = \exp(-n\log(1/q))$ dominates the subexponential
growth for any $q < 1$:
$q^n p(n) \leq C\exp(\alpha\sqrt{n} - n\log(1/q))$,
which decays superpolynomially.  This is the mechanism behind
the HS-sewing convergence for all standard families.
\end{remark}

\subsection{The Bergman space}\label{subsec:bergman}

Let $A^2(D)$ be the Bergman space of square-integrable
holomorphic functions on the unit disk
$D \subset \bC$, with inner product
$\langle f, g \rangle = \int_D f\overline{g}\,dA$ and
orthonormal basis
$e_n(z) = \sqrt{(n+1)/\pi}\,z^n$ for $n \geq 0$.
The symmetric algebra $\Sym A^2(D)$ is the bosonic Fock space
built on the one-particle Bergman space.  The Bergman
reproducing kernel is
\begin{equation}\label{eq:bergman-kernel}
K_D(z,w)
\;=\;
\frac{1}{\pi(1 - z\bar w)^2}.
\end{equation}

\begin{definition}[Mode--Bergman map]\label{def:theta-map}
The \emph{mode--Bergman map}
$\Theta \colon \cH_k^{\mathrm{alg}} \to \Sym A^2(D)$
is the dense linear map defined on the algebraic Fock core by
\begin{align*}
\Theta(a_{-n-1}\mathbf{1})
&\;=\;
\sqrt{n+1}\, z^n,\\
\Theta(a_{-n_1-1}\cdots a_{-n_j-1}\mathbf{1})
&\;=\;
\operatorname{Sym}\bigl(\Theta(a_{-n_1-1}\mathbf{1})
\otimes \cdots \otimes
\Theta(a_{-n_j-1}\mathbf{1})\bigr).
\end{align*}
\end{definition}

The map~$\Theta$ intertwines the algebraic Fock space grading
with the polynomial-degree grading of the Bergman space and
sends the invariant bilinear form on~$\cH_k$ to the Bergman
inner product (up to the level factor~$k$).

\subsection{The Heisenberg sewing theorem}
\label{subsec:heisenberg-sewing-theorem}

\begin{theorem}[Heisenberg sewing]
\label{thm:heisenberg-sewing}
Let $\cH_k$ be the Heisenberg vertex algebra at level~$k$ and
let~$\Theta$ be the mode--Bergman map.
\begin{enumerate}[label=\textup{(\roman*)}]
\item The sewing envelope~$\cH_k^{\mathrm{sew}}$ is
identified via~$\Theta$ with~$\Sym A^2(D)$.
\item The completed bar differential on Bergman tensors is the
closure of the Gaussian collision operator.
\item The pair-of-pants sewing amplitude is the second
quantisation~$\GG(R)$ of a one-particle Bergman restriction
$R \colon A^2(D_{\mathrm{out}}) \to
A^2(D_{\mathrm{in},1}) \otimes A^2(D_{\mathrm{in},2})$.
\item Every closed genus-$g$ amplitude is a Fredholm
determinant on the reduced one-particle Bergman space
$A^2_0(\partial D)$ (constant mode removed).
\end{enumerate}
\end{theorem}

\begin{proof}
Claim~(i) is the ind-Hilbert identification of
Moriwaki~\cite{Moriwaki26b}, who constructs a conformally flat
2-disk algebra structure on $\Sym A^2(D)$ isomorphic to the
Heisenberg sewing completion by factorisation homology valued
in $\operatorname{IndHilb}$.  The universal property
(Proposition~\ref{prop:sewing-universal}) identifies his
completion with~$\cH_k^{\mathrm{sew}}$.

Claim~(ii) follows from the analytic-algebraic comparison: the
algebraic bar differential, in the monomial basis of the Fock
space, is
\[
d_{\mathrm{bar}}
(a_{-n-1}\mathbf{1} \otimes a_{-m-1}\mathbf{1})
\;=\;
k\,\delta_{n+m+2,\,0}\,\mathbf{1}
\]
(the simple-pole Wick contraction of the generating current).
Under~$\Theta$, this contraction is the Gaussian collision
operator $G(z,w) = k/(1-z\bar w)^2$ on
$A^2(D) \otimes A^2(D)$, the reproducing kernel of the
Bergman space.

Claims~(iii) and~(iv) follow from the one-particle sewing
theorem below.
\end{proof}

\subsection{One-particle sewing and the Fredholm
 determinant}\label{subsec:one-particle}

The Heisenberg amplitude factorises completely through a
one-particle operator.  This is the analytic avatar of
Wick's theorem.

\begin{theorem}[One-particle sewing]
\label{thm:heisenberg-one-particle}
The Heisenberg pair-of-pants amplitude is the second
quantisation~$\GG(R)$ of the one-particle Bergman restriction
$R \colon A^2(D_{\mathrm{out}}) \to
A^2(D_{\mathrm{in},1}) \otimes A^2(D_{\mathrm{in},2})$,
with matrix elements satisfying
\[
|R_{n,m}| \;\leq\; C q^{(n+m)/2},
\]
where $q = e^{-t}$ and~$t$ is the conformal length of the
collar.  The sewing operator~$T_q$ on the reduced Bergman space
$A^2_0(\partial D)$ has eigenvalues~$q^n$ for $n \geq 1$
and is trace class with
\begin{equation}\label{eq:sewing-trace-norm}
\|T_q\|_1 \;=\; \frac{q}{1 - q}.
\end{equation}
The genus-$g$ partition function of~$\cH_k$ is the Fredholm
determinant
\begin{equation}\label{eq:heisenberg-fredholm}
Z_g(\cH_k;\,\Omega)
\;=\;
(\det\operatorname{Im}\Omega)^{-k/2}
\cdot (\det{}'_\zeta \Delta_{\Sigma_g})^{-k/2}.
\end{equation}
At genus~$1$,
$Z_1 = (\operatorname{Im}\tau)^{-k/2}\,
|\eta(\tau)|^{-2k}$.
\end{theorem}

\begin{proof}
Wick's theorem: every Heisenberg amplitude factors into
pairwise contractions of the generating current.  Each
contraction across a sewing circle corresponds to a
one-particle pairing, so the full Fock-space amplitude
equals the second quantisation~$\GG(R)$ of the one-particle
pairing~$R$.

The one-particle matrix element is
\[
R_{n,m}
\;=\;
\frac{k}{2\pi i}\oint_{|z| = 1}
\oint_{|w| = 1}
z^{-n-1}\, w^{-m-1}\,
G_{\mathrm{pants}}(z, w)\, dz\, dw,
\]
where $G_{\mathrm{pants}}$ is the Bergman kernel of the pair
of pants.  The collar of length $t = -\log q$ and the Bergman
reproducing property give
$|G_{\mathrm{pants}}(z, w)| \leq C(1 - q|zw|)^{-2}$;
extracting the $(n,m)$ Fourier coefficient yields
$|R_{n,m}| \leq C q^{(n+m)/2}$.

The operator~$R$ is Hilbert--Schmidt, so the
second-quantised operator~$\GG(R)$ is trace class, and
$\|\GG(R)\|_1 = \det(1 + |R|)^k < \infty$.

The Polyakov--Belavin--Knizhnik formula~\cite{BK86}
identifies the Heisenberg partition function:
$\det(1 - K_g)^{-k}
= (\det\operatorname{Im}\Omega)^{-k/2}
(\det{}'_\zeta \Delta_{\Sigma_g})^{-k/2}$,
where~$K_g$ is the genus-$g$ sewing kernel assembled by
composing $2g - 2$ pair-of-pants restriction operators along
the edges of a trivalent graph~$\GG$ of first Betti number~$g$.

At genus~$1$, $K_1 = T_q$ and the Fredholm determinant is
$\det(1 - T_q) = \prod_{n \geq 1}(1 - q^n) =
q^{-1/24}\eta(q)$,
so $Z_1 = (\operatorname{Im}\tau)^{-k/2}
\cdot |\eta(\tau)|^{-2k}$
after combining holomorphic and antiholomorphic sectors.
\end{proof}

\begin{remark}[Genus-1 connected free energy]
\label{rem:genus-1-free-energy}
At genus~$1$, the connected free energy is
\[
F_1^{\mathrm{conn}}(\cH_k;\,q)
\;=\;
-k\log\det(1 - T_q)
\;=\;
-k\log\bigl(q^{-1/24}\eta(q)\bigr)
\;=\;
k\Bigl(\frac{1}{24}\log q
- \sum_{n \geq 1}\log(1 - q^n)\Bigr),
\]
the Fourier expansion of which produces the connected sewing
amplitudes $a_{\cH_k}(N) = k\sum_{d \mid N} 1/d$; the
Dirichlet series is
$S_{\cH_k}(u) = k\,\zeta(u)\,\zeta(u+1)$.  The Gram
positivity of Corollary~\ref{cor:hs-positivity} is a
consequence of $a_{\cH_k}(N) > 0$.
\end{remark}

\subsection{Higher-genus Fredholm expansion}
\label{subsec:higher-genus-heisenberg}

The genus-$g$ partition function is assembled by choosing a
pants decomposition of~$\Sigma_g$: a trivalent graph~$\GG$ with
$g$ loops (first Betti number), $2g - 2$ vertices (pairs of
pants), and $3g - 3$ internal edges (sewing circles).  Each
vertex contributes a pair-of-pants restriction operator; each
edge contributes a sewing operator~$T_q$ (the restriction to a
collar of conformal length~$t = -\log q$).

The genus-$g$ sewing kernel is the composition
\begin{equation}\label{eq:genus-g-kernel}
K_g
\;=\;
\prod_{\text{edges of }\GG} T_{q_e}
\;\circ\;
\prod_{\text{vertices of }\GG} R_v,
\end{equation}
where $q_e = e^{-t_e}$ is the sewing parameter of the $e$-th
internal edge.  The eigenvalues of~$K_g$ on the one-particle
Bergman space are controlled by the Schottky parametrisation
of~$\Sigma_g$.

\begin{proposition}[Eigenvalue asymptotics of the sewing kernel]
\label{prop:sewing-eigenvalues}
The eigenvalues~$\lambda_j$ of the genus-$g$ sewing
kernel~$K_g$ satisfy:
\begin{enumerate}[label=\textup{(\roman*)}]
\item $0 \leq \lambda_j < 1$ for all~$j$;
\item $\sum_j \lambda_j < \infty$ (trace class);
\item $\sum_j \lambda_j^2 < \infty$ (Hilbert--Schmidt);
\item $\lambda_j \sim C_g j^{-1}\exp(-\alpha_g j^{1/g})$ as
 $j \to \infty$, where~$\alpha_g$ depends on the conformal
 structure of~$\Sigma_g$ through the extremal length of the
 separating geodesic.
\end{enumerate}
\end{proposition}

\begin{proof}
(i) and (ii): each~$T_{q_e}$ has eigenvalues $q_e^n$
($n \geq 1$), which are positive and summable.  Composition of
trace-class operators is trace class; eigenvalues of a
composition are bounded by the product of spectral radii.

(iii): the square-summability follows from the
Hilbert--Schmidt property of the pair-of-pants restriction
(Theorem~\ref{thm:heisenberg-one-particle}).

(iv): the asymptotic is determined by the extremal length of
the shortest separating geodesic on~$\Sigma_g$; the
exponential decay $\exp(-\alpha_g j^{1/g})$ reflects the
$g$-fold nature of the Schottky parametrisation.
\end{proof}

\begin{corollary}[Fredholm expansion at genus~$g$]
\label{cor:fredholm-expansion-g}
The genus-$g$ partition function admits the absolutely
convergent Fredholm expansion
\begin{equation}\label{eq:fredholm-expansion}
Z_g(\cH_k;\,\Omega)
\;=\;
\det(1 - K_g)^{-k}
\;=\;
\exp\Bigl(k\sum_{m \geq 1}
\frac{1}{m}\Tr(K_g^m)\Bigr),
\end{equation}
where the power traces $\Tr(K_g^m)$ are convergent sums over
products of sewing eigenvalues.
\end{corollary}

\begin{proof}
The Fredholm determinant of a trace-class operator satisfies
$\log\det(1 - K) = -\sum_{m \geq 1} \frac{1}{m}\Tr(K^m)$,
convergent whenever $\|K\|_1 < 1$.  The trace-class property
of~$K_g$ (Proposition~\ref{prop:sewing-eigenvalues}) ensures
convergence; raising to the $k$-th power accounts for the level
factor.
\end{proof}

\subsection{Genus-2 specialisation}\label{subsec:genus-2}

At genus~$2$, every Riemann surface admits a hyperelliptic
representation $y^2 = \prod_{i=1}^6(x - e_i)$, and the period
matrix is a $2 \times 2$ symmetric matrix
$\Omega \in \fH_2$ with $\operatorname{Im}\Omega > 0$.  The
Heisenberg partition function becomes
\begin{equation}\label{eq:genus-2-heisenberg}
Z_2(\cH_k;\,\Omega)
\;=\;
(\det\operatorname{Im}\Omega)^{-k/2}
\cdot (\det{}'_\zeta \Delta_{\Sigma_2})^{-k/2}.
\end{equation}
The $2 \times 2$ determinant
$\det\operatorname{Im}\Omega
= (\operatorname{Im}\Omega_{11})
(\operatorname{Im}\Omega_{22})
- (\operatorname{Im}\Omega_{12})^2$
is the volume of the Jacobian torus
$\bC^2/(\bZ^2 + \Omega\bZ^2)$.

The sewing kernel at genus~$2$ is assembled from a trivalent
graph with $2$ vertices and $3$ internal edges:
$K_2 = T_{q_1}\circ R_1\circ T_{q_2}\circ R_2\circ T_{q_3}$,
where $q_1, q_2, q_3$ are the three sewing parameters
corresponding to the three handles of the genus-$2$ surface
(two non-separating, one separating).  The degeneration limits
$q_i \to 0$ recover the genus-$1$ partition functions of the
constituent tori.

\begin{remark}[Factorisation at the separating
 node]\label{rem:separating-node}
In the separating degeneration (one of the three sewing
parameters approaches zero while the other two remain generic),
the genus-$2$ partition function factors:
\[
Z_2(\cH_k;\,\Omega)
\;\xrightarrow{q_3 \to 0}\;
Z_1(\cH_k;\,\tau_1) \cdot Z_1(\cH_k;\,\tau_2) \cdot
(1 + O(q_3)),
\]
where $\tau_1, \tau_2$ are the modular parameters of the two
genus-$1$ components.  The leading term is the product of
two genus-$1$ partition functions; the corrections are
trace-class contributions from the sewing circle.  The
non-separating degeneration carries genuinely genus-$2$
information: the $2 \times 2$ period matrix controls the
interplay of the two homology cycles.
\end{remark}

\begin{remark}[The Heisenberg as the atom of sewing]
\label{rem:heisenberg-atom}
The Heisenberg algebra is the unique member of the standard
landscape for which every step of the sewing programme is
explicit: the sewing envelope is a classical Bergman space;
the pair-of-pants amplitude is a Fock-space second
quantisation; every partition function is a Fredholm
determinant.  All non-abelian chiral algebras satisfying
HS-sewing require the abstract approach of
Section~\ref{sec:general-criterion}.  The lattice vertex
algebra is the first extension (Section~\ref{sec:lattice}),
where the Heisenberg Fredholm determinant appears as one
factor of the full partition function.
\end{remark}


\subsection{The Heisenberg at higher genus: Schottky
 uniformisation}\label{subsec:schottky}

The Fredholm determinant formula for the Heisenberg partition
function admits an interpretation in terms of Schottky
uniformisation that illuminates the sewing structure.

A genus-$g$ Schottky group is generated by $g$
loxodromic M\"obius transformations
$\gamma_1,\ldots,\gamma_g \in \mathrm{PSL}_2(\bC)$, each
with multiplier~$q_j$ ($|q_j| < 1$).  The Schottky
uniformisation represents $\Sigma_g$ as a quotient of a
domain in~$\bC\mathrm{P}^1$ by the Schottky group.

The sewing construction realises this uniformisation at the
chain level: each generator~$\gamma_j$ corresponds to a sewing
circle with parameter~$q_j$, and the Fredholm determinant is
\begin{equation}\label{eq:schottky-fredholm}
Z_g(\cH_k;\,q_1,\ldots,q_g)
\;=\;
\prod_{n \geq 1}
\prod_{\{\gamma\}'} (1 - q_\gamma^n)^{-k},
\end{equation}
where $\{\gamma\}'$ denotes the set of primitive conjugacy
classes in the Schottky group and $q_\gamma$ is the multiplier
of~$\gamma$.  The product converges absolutely because the
multipliers decay exponentially: $|q_\gamma| \leq
\max_j|q_j|^{|\gamma|}$, where $|\gamma|$ is the word length.

\begin{remark}[Selberg zeta function]
\label{rem:selberg-zeta}
The Schottky product~\eqref{eq:schottky-fredholm} is the
inverse of the Selberg zeta function
$Z_{\mathrm{Selberg}}(\Sigma_g,s)$ evaluated at $s = 1$.
The Selberg zeta function encodes the spectrum of the
Laplacian on~$\Sigma_g$; its value at $s = 1$ is
$\det{}'_\zeta\Delta_{\Sigma_g}$.  The Heisenberg partition
function is therefore
$Z_g(\cH_k) = Z_{\mathrm{Selberg}}(\Sigma_g, 1)^{-k/2}$
times the period-matrix factor
$(\det\operatorname{Im}\Omega)^{-k/2}$.  This connects the
algebraic sewing construction to spectral geometry.
\end{remark}

\begin{remark}[The Polyakov conjecture]
\label{rem:polyakov-conjecture}
Polyakov's programme computes the string partition function as a
product of determinants: the holomorphic sector gives
$(\det{}'_\zeta\Delta)^{-k/2}$ (the zeta-regularised
determinant of the scalar Laplacian), and the non-holomorphic
sector gives $(\det\operatorname{Im}\Omega)^{-k/2}$ (the
volume of the Jacobian torus).  The HS-sewing theorem provides
the vertex-algebraic proof: the algebraic sewing
(Fredholm determinant of the one-particle Bergman operator)
reproduces the spectral-geometric answer (zeta-regularised
determinant of the Laplacian), and both are equal to the
Belavin--Knizhnik formula~\cite{BK86}.  The three independent
computations (algebraic, spectral, geometric) constitute a
verification of the Heisenberg sewing theorem by independent
paths.
\end{remark}


% ================================================================
% 5. LATTICE SEWING: THETA FUNCTIONS
% ================================================================

\section{Lattice sewing: theta functions}\label{sec:lattice}

The lattice vertex algebra~$V_\Lambda$ associated to a
positive-definite even lattice $\Lambda$ of rank~$r$ provides
the natural next case after the Heisenberg: the sewing
programme extends by combining the Heisenberg Fredholm
determinant with the convergence of the Siegel theta function.

\subsection{Structure of the lattice vertex algebra}
\label{subsec:lattice-structure}

Let~$\Lambda$ be a positive-definite even lattice of rank~$r$,
with dual lattice~$\Lambda^*$ and discriminant group
$D(\Lambda) = \Lambda^*/\Lambda$.  The lattice vertex algebra
$V_\Lambda$ is the tensor product of the rank-$r$ Heisenberg
Fock space with the group algebra of~$\Lambda$:
\[
V_\Lambda
\;=\;
\bigoplus_{\gamma \in D(\Lambda)} \cF_\gamma,
\qquad
\cF_\gamma \;=\; \cH_r \otimes \bC[[\Lambda + \gamma]],
\]
where~$\cF_\gamma$ is the twisted Heisenberg module indexed by
the coset~$\gamma \in D(\Lambda)$.  The vertex operators
$e^\gamma$ ($\gamma \in \Lambda$) act by
\[
e^\gamma \cdot
a_{-n_1}\cdots a_{-n_j} |\gamma'\rangle
\;=\;
\varepsilon(\gamma,\gamma') \cdot
a_{-n_1}\cdots a_{-n_j} |\gamma + \gamma'\rangle,
\]
where $\varepsilon \colon \Lambda \times \Lambda \to
\{\pm 1\}$ is the 2-cocycle.  The central charge is~$c = r$
and $\kappa(V_\Lambda) = r$.

Sector dimensions satisfy
$\dim(V_\Lambda)_n = |D(\Lambda)| \cdot p_r(n)$,
where~$p_r(n)$ is the $r$-coloured partition number.  The
sector growth is subexponential:
$\log\dim(V_\Lambda)_n = O(\sqrt{n})$.  OPE structure
constants satisfy
$|C^{c,k}_{a,i;\,b,j}| \leq K(a+b+c+1)^r$ because all OPE
singularities in~$V_\Lambda$ are at most $r$-th order poles.

\subsection{The lattice partition function}
\label{subsec:lattice-partition}

\begin{theorem}[Lattice sewing]\label{thm:lattice-sewing}
Let $\Lambda$ be a positive-definite even lattice of rank~$r$.
\begin{enumerate}[label=\textup{(\roman*)}]
\item (\textbf{HS-sewing})\; $V_\Lambda$ satisfies HS-sewing
 at every $0 < q < 1$.
\item (\textbf{Factorisation})\; The genus-$g$ partition
 function factors as
 \begin{equation}\label{eq:lattice-partition}
 Z_g(V_\Lambda;\,\Omega)
 \;=\;
 \underbrace{(\det\operatorname{Im}\Omega)^{-r/2}
 \cdot (\det{}'_\zeta \Delta_{\Sigma_g})^{-r/2}}_{%
  Z_g(\cH_r;\,\Omega)}
 \;\cdot\;
 \underbrace{\sum_{\lambda \in \Lambda^g}
  e^{i\pi \lambda^T \Omega \lambda}}_{%
  \Theta_\Lambda(\Omega)},
 \end{equation}
 and each factor converges absolutely on the Siegel upper
 half-space~$\fH_g$.
\item (\textbf{Sewing envelope})\; The sewing envelope is the
 charge-refined Hilbert completion
 \begin{equation}\label{eq:lattice-sewing-envelope}
 V_\Lambda^{\mathrm{sew}}
 \;=\;
 \widehat{\bigoplus}_{\gamma \in D(\Lambda)}
 \Sym\, A^2_\gamma(D),
 \end{equation}
 where $A^2_\gamma(D) \subset A^2(D)$ is the Bergman space
 with~$L_0$-shift
 $\tfrac{1}{2}\langle\gamma,\gamma\rangle$ and the completion
 is in the HS-sewing topology.
\item (\textbf{Bounded vertex operators})\; The vertex
 operators $e^\gamma \colon
 \widehat{\cF}_{\gamma'} \to
 \widehat{\cF}_{\gamma'+\gamma}$ extend to bounded operators
 between completed sectors with
 $\|e^\gamma\|_{\mathrm{op}}
 \leq q^{\langle\gamma,\gamma\rangle/2}$.
\end{enumerate}
\end{theorem}

\begin{proof}
(i) follows from
$\log\dim(V_\Lambda)_n = O(\sqrt{n})$ (subexponential sector
growth) and $|C^{c,k}_{a,i;\,b,j}| \leq K(a+b+c+1)^r$
(polynomial OPE growth): the hypotheses of
Theorem~\ref{thm:general-hs-sewing} are satisfied.

(ii): the Fock decomposition
$V_\Lambda = \bigoplus_\gamma \cF_\gamma$ is a direct sum
of twisted Heisenberg modules.  Within each sector, Wick's
theorem reduces all correlation functions to propagator
pairings on the underlying rank-$r$ Heisenberg, contributing
$Z_g(\cH_r;\,\Omega)$
(Theorem~\ref{thm:heisenberg-one-particle}).  The lattice sum
over winding modes contributes~$\Theta_\Lambda(\Omega)$.  The
Heisenberg factor converges by the one-particle sewing theorem
(eigenvalues~$q^n$, trace class).  The theta factor converges
absolutely on~$\fH_g$: the defining sum satisfies
$|\Theta_\Lambda(\Omega)|
\leq \sum_\lambda e^{-\pi\lambda^T Y\lambda}$
with $Y = \operatorname{Im}(\Omega) > 0$, and the exponent
$\lambda^T Y\lambda
\geq \lambda_{\min}(Y) \cdot \|\lambda\|^2$
grows as $\|\lambda\|^2 \to \infty$.

(iii): the Bergman-space completion
$\widehat{\cF}_\gamma
= \overline{\cF_\gamma}^{\|\cdot\|_{A^2}}$
gives honest Hilbert sectors.  The direct-sum decomposition of
Proposition~\ref{prop:sewing-direct-sum} applies to the
$D(\Lambda)$-grading.

(iv): the vertex operator~$e^\gamma$ preserves mode operators
and only shifts the lattice vector.  The cocycle~$\varepsilon$
is bounded by~$1$.  The $L_0$-shift
$\Delta_\gamma = \frac{1}{2}\langle\gamma,\gamma\rangle$ is
absorbed into the grading; the exponential weight damping
$q^{\Delta_\gamma}$ in the HS topology gives the bound.
\end{proof}

\begin{remark}[Self-dual lattices and modular
 invariance]\label{rem:self-dual}
When $\Lambda$ is unimodular
($D(\Lambda) = \{0\}$, e.g.\ $E_8$ or
$\Lambda_{24}$), the theta function
$\Theta_\Lambda(\Omega)$ is a Siegel modular form with
respect to $\mathrm{Sp}_{2g}(\bZ)$.  The full partition
function $Z_g(V_\Lambda;\,\Omega)$ then transforms as a
modular form of weight~$r/2$, recovering Zhu's modular
invariance theorem~\cite{Zhu96} from the sewing formalism.
For non-self-dual lattices, the theta function transforms
under a finite-index congruence subgroup of
$\mathrm{Sp}_{2g}(\bZ)$, and the partition function has the
corresponding transformation law.
\end{remark}

\begin{remark}[Genus-1 specialisation]
\label{rem:lattice-genus-1}
At genus~$1$, the lattice partition function reduces to
\[
Z_1(V_\Lambda;\,\tau)
\;=\;
\frac{\Theta_\Lambda(\tau)}{|\eta(\tau)|^{2r}}.
\]
For the $E_8$ lattice ($r = 8$): the theta function
$\Theta_{E_8}(\tau) = E_4(\tau)$ (the Eisenstein series of
weight~$4$), and
$Z_1(V_{E_8};\,\tau)
= E_4(\tau)/|\eta(\tau)|^{16}$.
For the Leech lattice ($r = 24$):
$\Theta_{\Lambda_{24}}(\tau) = E_{12}(\tau) + \cdots$
(a weight-$12$ modular form).
\end{remark}

\begin{example}[Genus-2 for the $E_8$ lattice]
\label{ex:e8-genus2}
At genus~$2$ with period matrix
$\Omega = \begin{pmatrix}\tau_1 & z \\ z & \tau_2
\end{pmatrix} \in \fH_2$, the $E_8$ lattice partition function
is
\[
Z_2(V_{E_8};\,\Omega)
\;=\;
Z_2(\cH_8;\,\Omega) \cdot \Theta_{E_8}(\Omega),
\]
where the Heisenberg factor is
$Z_2(\cH_8;\,\Omega) = (\det\operatorname{Im}\Omega)^{-4}
\cdot (\det{}'_\zeta\Delta_{\Sigma_2})^{-4}$
and the genus-$2$ theta function is
\[
\Theta_{E_8}(\Omega)
\;=\;
\sum_{\lambda \in E_8^2}
e^{i\pi\lambda^T\Omega\lambda}
\;=\;
\sum_{\lambda_1,\lambda_2 \in E_8}
e^{i\pi(\tau_1|\lambda_1|^2 + 2z\langle\lambda_1,
\lambda_2\rangle + \tau_2|\lambda_2|^2)}.
\]
The genus-$2$ theta function
$\Theta_{E_8}(\Omega)$ is a Siegel modular form of weight~$4$
with respect to $\mathrm{Sp}_4(\bZ)$.  The factorisation of
$Z_2$ into Heisenberg and lattice parts extends the genus-$1$
formula $Z_1 = E_4(\tau)/|\eta(\tau)|^{16}$ to genus~$2$.

The convergence of~$\Theta_{E_8}(\Omega)$ on~$\fH_2$ is
controlled by the shortest vector of~$E_8$
($|\lambda_{\min}|^2 = 2$): the tail of the theta sum is
bounded by
$\sum_{|\lambda| \geq R}
e^{-\pi\lambda_{\min}(Y)\cdot|\lambda|^2}$,
which decays as $e^{-\pi\lambda_{\min}(Y)\cdot R^2}$.
\end{example}

\begin{example}[The $D_{16}^+$ lattice and moonshine]
\label{ex:d16plus}
The lattice $D_{16}^+$ (rank~$16$, even, unimodular, but not
isomorphic to $E_8 \oplus E_8$) provides a lattice VOA with
the same genus-$1$ partition function as
$V_{E_8 \oplus E_8}$ (both equal
$E_4(\tau)^2/|\eta(\tau)|^{32}$).  The two VOAs are
distinguished at genus~$2$: their Siegel theta functions
$\Theta_{D_{16}^+}(\Omega) \neq
\Theta_{E_8 \oplus E_8}(\Omega)$ are distinct Siegel modular
forms of weight~$8$.  The analytic sewing programme detects
this distinction through the genus-$2$ partition function,
providing finer invariants than the genus-$1$ character.
\end{example}

\subsection{Convergence of the Siegel theta
 function}\label{subsec:theta-convergence}

The theta factor $\Theta_\Lambda(\Omega)
= \sum_{\lambda \in \Lambda^g}
e^{i\pi\lambda^T\Omega\lambda}$ in the lattice partition
function requires a separate convergence analysis from the
Heisenberg Fredholm factor.  The two convergence mechanisms are
logically independent.

\begin{proposition}[Theta convergence on the Siegel upper
 half-space]\label{prop:theta-convergence}
Let $\Lambda$ be a positive-definite even lattice of rank~$r$.
For any $\Omega \in \fH_g$
(i.e.\ $\Omega^T = \Omega$ and
$Y := \operatorname{Im}\Omega > 0$), the theta series
converges absolutely with the bound
\begin{equation}\label{eq:theta-bound}
|\Theta_\Lambda(\Omega)|
\;\leq\;
\sum_{\lambda \in \Lambda^g}
e^{-\pi\lambda^T Y\lambda}
\;\leq\;
C(Y)^{r/2}
\cdot \bigl(\det Y\bigr)^{-r/2},
\end{equation}
where $C(Y)$ is a constant depending on the lattice through the
shortest vector.
\end{proposition}

\begin{proof}
For $\lambda = (\lambda_1,\ldots,\lambda_g) \in \Lambda^g$,
the real part of $i\pi\lambda^T\Omega\lambda$ is
$-\pi\lambda^T Y\lambda$.  Since $Y > 0$, each summand satisfies
$|e^{i\pi\lambda^T\Omega\lambda}|
= e^{-\pi\lambda^T Y\lambda}$.
The quadratic form
$\lambda^T Y\lambda
\geq \lambda_{\min}(Y)
\cdot \|\lambda\|^2$
grows as $\|\lambda\|^2$, so the sum converges by comparison
with a Gaussian integral over~$\bR^{rg}$.  The bound follows
from the Poisson summation formula applied to the Gaussian
$e^{-\pi\lambda^T Y\lambda}$.
\end{proof}

\begin{remark}[Two independent convergence
 mechanisms]\label{rem:two-mechanisms}
The lattice partition function~\eqref{eq:lattice-partition}
has two factors, each with a distinct convergence mechanism:
\begin{enumerate}[label=\textup{(\alph*)}]
\item the Heisenberg factor~$Z_g(\cH_r;\,\Omega)$ converges by
 the Fredholm determinant (trace-class sewing operator,
 eigenvalue decay $q^n$, $n \geq 1$);
\item the theta factor~$\Theta_\Lambda(\Omega)$ converges by
 positive-definiteness of~$\operatorname{Im}\Omega$
 (Gaussian decay of the lattice sum).
\end{enumerate}
Neither mechanism alone suffices for the other factor.  The
Heisenberg factor involves the Bergman kernel and the Arakelov
Green function; the theta factor involves the lattice geometry
and the period matrix.  The factorisation
$Z_g = Z_g^{\mathrm{Heis}} \cdot \Theta_\Lambda$ is a
consequence of the Fock decomposition of~$V_\Lambda$ into
twisted Heisenberg modules; no analogous factorisation is
available for non-abelian algebras.
\end{remark}

\subsection{The Dirichlet--sewing lift for
 lattice VOAs}\label{subsec:lattice-dirichlet}

The connected free energy of the lattice VOA at genus~$1$ is
\begin{equation}\label{eq:lattice-f1}
F_1^{\mathrm{conn}}(V_\Lambda;\,q)
\;=\;
-r\log\det(1 - T_q)
- \log\Theta_\Lambda(\tau),
\end{equation}
where the first term is the rank-$r$ Heisenberg contribution
and the second is the lattice correction.  The Dirichlet--sewing
lift produces
\begin{equation}\label{eq:lattice-dirichlet}
S_{V_\Lambda}(u)
\;=\;
r\,\zeta(u)\,\zeta(u+1)
+ S_\Lambda^{\mathrm{lattice}}(u),
\end{equation}
where $S_\Lambda^{\mathrm{lattice}}(u)$ encodes the theta
function's Fourier coefficients as a Dirichlet series.  For
self-dual lattices, the lattice Dirichlet series inherits the
modular invariance of the theta function.


% ================================================================
% 6. GENERAL HS-SEWING THEOREM
% ================================================================

\section{General HS-sewing theorem for the standard landscape}
\label{sec:general-criterion}

For chiral algebras beyond the Heisenberg and lattice families,
the sewing envelope does not admit a classical description, but
the Hilbert--Schmidt criterion can still be verified from two
OPE-level estimates.

\subsection{The general criterion}\label{subsec:general-criterion}

\begin{theorem}[General HS-sewing criterion]
\label{thm:general-hs-sewing}
Let $\cA$ be a positive-energy chiral algebra with a choice of
pre-Hilbert structure on each sector~$H_n$.  Suppose:
\begin{enumerate}[label=\textup{(\roman*)}]
\item \textup{(\textbf{subexponential sector growth})}\;
$\log\dim H_n = o(n)$; equivalently,
$\dim H_n \leq C e^{\alpha\sqrt{n}}$ for some
$C, \alpha > 0$;
\item \textup{(\textbf{polynomial OPE coefficient growth})}\;
there exist $K > 0$ and $N \geq 0$ such that the matrix
elements of~$m_{a,b}^c$ in orthonormal bases satisfy
\[
|C^{c,k}_{a,i;\,b,j}|
\;\leq\;
K(a + b + c + 1)^N
\]
for all weights $a,b,c$ and all basis indices $i,j,k$.
\end{enumerate}
Then $\cA$ satisfies HS-sewing at every parameter
$0 < q < 1$.
\end{theorem}

\begin{proof}
From hypothesis~(ii), the Hilbert--Schmidt norm of the sector
multiplication is bounded by
\begin{equation}\label{eq:hs-bound}
\|m_{a,b}^c\|_{\HS}^2
\;\leq\;
\dim H_a \cdot \dim H_b \cdot \dim H_c \cdot
K^2(a+b+c+1)^{2N}.
\end{equation}
The power-mean inequality
$(x+y+z)^p \leq 3^{p-1}(x^p + y^p + z^p)$ with $p = 2N$ gives
\[
(a+b+c+1)^{2N}
\;\leq\;
3^{2N-1}\bigl((a+1)^{2N} + (b+1)^{2N} + (c+1)^{2N}\bigr).
\]
Substituting into~\eqref{eq:hs-bound} and using the symmetry
of the three terms:
\[
\sum_{a,b,c \geq 0}
q^{a+b+c}\,\|m_{a,b}^c\|_{\HS}^2
\;\leq\;
3^{2N}K^2
\Bigl(\sum_{n \geq 0} q^n\dim H_n\,(n+1)^{2N}\Bigr)
\Bigl(\sum_{n \geq 0} q^n\dim H_n\Bigr)^2.
\]
By hypothesis~(i), $\dim H_n \leq Ce^{\alpha\sqrt{n}}$,
while $q^n = e^{-n\log(1/q)}$ decays exponentially.  Each
factor $q^n\dim H_n\,(n+1)^{2N}$ is bounded by
$C(n+1)^{2N}\exp(\alpha\sqrt{n} - n\log(1/q))$,
which decays superpolynomially.  Each series converges
absolutely.
\end{proof}

\begin{corollary}[Automatic HS-sewing for Koszul
 algebras]\label{cor:koszul-automatic}
Every finitely strongly generated chirally Koszul algebra
satisfies HS-sewing.  The PBW theorem provides polynomial OPE
growth; finite strong generation provides the finite $N$ in
hypothesis~(ii); positive energy provides subexponential sector
growth.  HS-sewing is therefore an independent analytic
condition only for infinitely generated algebras.
\end{corollary}

\subsection{Verification on the standard
 landscape}\label{subsec:verification}

We verify the hypotheses of
Theorem~\ref{thm:general-hs-sewing} for each standard family.
The data required is a character formula (controlling
$\dim H_n$) and a polynomial OPE bound (determining~$N$).

\begin{example}[Heisenberg]\label{ex:heisenberg}
The Heisenberg~$\cH_k$ has $\dim H_n = p(n)$, the partition
number.  The OPE
$J(z)J(w) \sim k/(z-w)^2$ has structure constants bounded
by~$k$, independent of weight; $N = 1$ suffices (the
derivation~$\partial$ increases weight by~$1$, so normal-ordered
products of derivatives are at most linear in the weight).
\end{example}

\begin{example}[Affine Kac--Moody]\label{ex:kac-moody}
Let $V_k(\fg)$ be the vacuum module at level~$k$, with $\fg$
simple.  Sectors satisfy
$\dim H_n \leq Ce^{\alpha\sqrt{n}}$ from the Kac--Weyl
character formula.  The OPE
$J^a(z)J^b(w) \sim k\,\delta^{ab}/(z-w)^2
+ f^{ab}{}_c J^c(w)/(z-w)$
gives $|C| \leq K(n+1)$ with $N = 1$.  HS-sewing at every
$0 < q < 1$ follows.
\end{example}

\begin{example}[Virasoro]\label{ex:virasoro}
The Virasoro algebra~$\Vir_c$ has $\dim H_n = p(n)$.  The OPE
$T(z)T(w) \sim (c/2)(z-w)^{-4} + 2T(w)(z-w)^{-2}
+ \partial T(w)(z-w)^{-1}$
is quartic (the $r$-matrix has cubic pole by $d\log$-absorption,
per the pole relation $\mathrm{pole}_r = \mathrm{pole}_{\mathrm{OPE}} - 1$).
Repeated normal products give $|C| \leq K(n+1)^2$ with
$N = 2$.  The criterion applies at every central charge~$c$.
\end{example}

\begin{example}[$\cW_N$-algebras]\label{ex:wn}
The principal~$\cW_N$-algebra has $N-1$ generators of
conformal weights $2,3,\ldots,N$.  Sector dimensions satisfy
$\dim H_n \leq Ce^{\alpha\sqrt{n}}$ from the $\cW_N$ character
formula.  OPE coefficient growth is
$|C| \leq K(n+1)^{2N}$; applying
Theorem~\ref{thm:general-hs-sewing} gives HS-sewing at every
$0 < q < 1$.
\end{example}

\begin{example}[Lattice vertex algebras]\label{ex:lattice}
Covered by Theorem~\ref{thm:lattice-sewing}; recorded here
for completeness.  $\dim H_n = |D(\Lambda)| \cdot p_r(n)$,
$|C| \leq K \cdot r$, $N = r$.
\end{example}

\begin{proposition}[Summary table]\label{prop:summary-table}
The following table records the HS-sewing data for the standard
landscape.  Every entry satisfies the hypotheses of
Theorem~\ref{thm:general-hs-sewing}.
\begin{center}
\renewcommand{\arraystretch}{1.15}
\begin{tabular}{l|ccc|c}
\toprule
$\cA$ & $\dim H_n$ & $|C|$ bound & $N$ &
 $\sum_n q^n\dim H_n$ \\
\midrule
$\cH_k$ & $p(n)$ & $\leq k$
 & $1$ & $q^{1/24}/\eta(q)$ \\
$V_k(\fg)$ & $\leq Ce^{\alpha\sqrt{n}}$
 & $\leq K(n+1)$ & $1$ & $\chi_{V_k}(q)$ \\
$\Vir_c$ & $p(n)$
 & $\leq K(n+1)^2$ & $2$
 & $q^{(1-c)/24}/\eta(q)$ \\
$\cW_N$ & $\leq Ce^{\alpha\sqrt{n}}$
 & $\leq K(n+1)^{2N}$ & $2N$
 & $\chi_{\cW_N}(q)$ \\
$V_\Lambda$ & $|D(\Lambda)|\cdot p_r(n)$
 & $\leq Kr$ & $r$
 & $\Theta_\Lambda(q)/\eta(q)^r$ \\
\bottomrule
\end{tabular}
\end{center}
\end{proposition}

\begin{remark}[Completed algebras]\label{rem:completed}
The standard families above are finitely strongly generated.
For completed algebras (such as $\cW_{1+\infty}$ or completed
Yangians), the polynomial OPE bound of hypothesis~(ii) must be
verified in the completed (inverse-limit) topology, where the
structure constants are formal power series.  This verification
is the content of the MC4 completion programme
of~\cite{LorgatVolI}.
\end{remark}

\begin{remark}[Beyond the standard landscape]
\label{rem:beyond-standard}
The HS-sewing criterion is a priori applicable to any
positive-energy chiral algebra satisfying subexponential growth
and polynomial OPE bounds.  Two classes of algebras beyond the
standard landscape merit investigation:
\begin{enumerate}[label=\textup{(\alph*)}]
\item Admissible-level simple quotients~$L_k(\fg)$ satisfy
 $C_2$-cofiniteness~\cite{Zhu96} and hence HS-sewing (the
 $C_2$-cofinite condition implies convergent characters, which
 implies subexponential growth; the simple quotient has
 finite-dimensional fused products, which implies polynomial
 OPE bounds).
\item $\cW$-algebras at non-generic central charge, where the
 sector growth may exhibit resonance phenomena but remains
 subexponential.
\end{enumerate}
\end{remark}

\subsection{Fredholm expansion and the shadow
 tower}\label{subsec:fredholm-shadow}

The HS-sewing theorem delivers convergent partition functions;
the Fredholm expansion connects them to the shadow obstruction
tower of the modular Koszul duality programme.

For any HS-sewing algebra, the genus-$g$ partition function
admits a graph-sum expansion over stable graphs
$\GG \in \cS_g$ (the set of stable graphs
of~$\overline\cM_g$):
\begin{equation}\label{eq:graph-sum}
Z_g(\cA;\,\Omega)
\;=\;
\sum_{\GG \in \cS_g}
\frac{1}{|\mathrm{Aut}(\GG)|}
\prod_{v \in V(\GG)} A_v(\cA)
\prod_{e \in E(\GG)} K_e(q_e),
\end{equation}
where $A_v$ is the vertex amplitude (genus-$g_v$ with $n_v$
legs) and $K_e(q_e)$ is the sewing kernel along edge~$e$.
Each term is absolutely convergent by the trace-class property.

The shadow obstruction tower
$\Theta_\cA = \sum_{r \geq 2} S_r(\cA)$
governs the leading-order asymptotics of the graph-sum.
At uniform weight, the genus-$g$ partition function satisfies
$F_g(\cA) = \kappa(\cA) \cdot \lambda_g^{\mathrm{FP}}$
\textsc{(uniform-weight)}, where $\lambda_g^{\mathrm{FP}}$ is
the Faber--Pandharipande evaluation of the tautological
class~$\lambda_g$ on~$\overline\cM_g$.

\begin{proposition}[Convergence of the shadow tower evaluation]
\label{prop:shadow-convergence}
Under HS-sewing, the evaluation of the MC element
$\Theta_\cA$ on any tautological class
$\alpha \in H^*(\overline\cM_{g,n})$ converges to a finite
complex number.  In particular, the genus-$g$ shadow
obstruction $F_g(\cA) = \langle\Theta_\cA, [\overline\cM_g]
\rangle$ is finite.
\end{proposition}

\begin{proof}
The MC element $\Theta_\cA = D_\cA - d_0$ lies in the modular
convolution algebra~$\mathfrak{g}^{\mathrm{mod}}_\cA$.  Its
evaluation on~$\alpha$ is a sum of graph amplitudes
$\langle\Theta_\cA^{(r)}, \alpha\rangle$ indexed by the
degree~$r$ of the MC component.  Each graph amplitude factors
through the trace-class sewing operators by the same mechanism
as the partition function; the trace-class property ensures
absolute convergence of each term and of the sum over~$r$.
\end{proof}

\subsection{The Dirichlet--sewing lift}\label{subsec:dirichlet-lift}

The connected free energy at genus~$1$ encodes connected
sewing amplitudes $a_\cA(N)$ as Fourier coefficients.  The
Dirichlet--sewing lift
\begin{equation}\label{eq:dirichlet-sewing}
S_\cA(u)
\;=\;
\sum_{N \geq 1} a_\cA(N) N^{-u}
\;=\;
\frac{(2\pi)^u}{\GG(u)}
\int_0^\infty F_\cA^{\mathrm{conn}}(e^{-2\pi y})\,
y^{u-1}\,dy
\end{equation}
is a Dirichlet series whose analytic structure is controlled by
the MC element~$\Theta_\cA$.

\begin{example}[Explicit Dirichlet series]
\label{ex:dirichlet-explicit}
The sewing Dirichlet series for the standard families:
\begin{align}
S_{\cH_k}(u) &= k\,\zeta(u)\,\zeta(u+1),
\label{eq:dirichlet-heisenberg}\\
S_{V_k(\fg)}(u) &= \dim(\fg)\cdot\zeta(u)\,\zeta(u+1),
\label{eq:dirichlet-km}\\
S_{\Vir_c}(u) &= \zeta(u+1)\bigl(\zeta(u) - 1\bigr),
\label{eq:dirichlet-virasoro}\\
S_{\cW_N}(u) &= \zeta(u+1)\Bigl((N-1)\zeta(u)
- \textstyle\sum_{j=1}^{N-1} H_j(u)\Bigr),
\label{eq:dirichlet-wn}
\end{align}
where $H_j(u) = \sum_{n=1}^j n^{-u}$ are partial harmonic
sums.  The Heisenberg and Kac--Moody series have a double pole
at $u = 0$; the Virasoro series has a simple pole (the
subtraction of~$1$ from $\zeta(u)$ removes the $N = 1$ term).
\end{example}

The sewing amplitudes $a_\cA(N)$ define a positive-semidefinite
inner product
$\langle f, g\rangle_\cA
= \sum_N a_\cA(N) f(N) g(N)$
with reproducing kernel
$K_\cA(s,t) = S_\cA(s+t)$.  The surface moment matrix
$\mathscr{M}_{ij}
= K_\cA(\alpha/2 + i, \alpha/2 + j)$
is positive definite for all $\alpha \geq 2$: a Gram matrix
$V^T D V$ with Vandermonde~$V$.


% ================================================================
% 7. THE THREE-LAYER ANALYTIC GAP
% ================================================================

\section{The three-layer analytic gap}\label{sec:gap}

The HS-sewing theorem guarantees convergent partition functions
for the standard landscape, but three layers of structure
separate the convergence statement from a complete analytic
theory.

\subsection{Layer 1: the sewing envelope for interacting
 algebras}\label{subsec:gap-layer1}

For the Heisenberg algebra, the sewing envelope is
$\Sym A^2(D)$: a classical Fock space built on the
Bergman space of the disk.  For lattice vertex algebras, it is
the charge-refined Hilbert
completion~\eqref{eq:lattice-sewing-envelope}.  Both
identifications rest on the same mechanism: the Wick theorem
reduces all amplitudes to one-particle (Heisenberg) or
one-particle-plus-lattice (lattice) data, and the sewing
envelope inherits the Bergman structure from the one-particle
space.

For non-abelian chiral algebras (affine Kac--Moody at generic
level, Virasoro, $\cW_N$), no such reduction is available.
The sewing envelope exists abstractly by
Definition~\ref{def:sewing-envelope}, and HS-sewing guarantees
that the weighted completion~$H_q$ is a concrete Hilbert space
on which all amplitudes are trace class.  The open problem is a
structural identification of~$\cA^{\mathrm{sew}}$ analogous
to the Bergman identification for the Heisenberg.

\begin{conjecture}[Sewing envelope for
 $V_k(\fg)$]\label{conj:sewing-envelope-km}
The sewing envelope of $V_k(\fg)$ at integer level
$k \in \bZ_{>0}$ is the ind-Hilbert completion
$\widehat{V_k(\fg)}^{\mathrm{IndHilb}}$, and the
pair-of-pants amplitude admits a Fredholm expansion indexed
by Weyl-orbit data.
\end{conjecture}

The next step is explicit: the case $\fg = \mathfrak{sl}_2$ at
$k = 1$ is the simplest interacting sewing computation.

\begin{example}[$\mathfrak{sl}_2$ at level~$1$]
\label{ex:sl2-level1}
The vacuum module $V_1(\mathfrak{sl}_2)$ has central charge
$c = 1$ and $\kappa = \dim(\mathfrak{sl}_2)(1+2)/(2\cdot 2)
= 3 \cdot 3/4 = 9/4$.  At conformal weight~$2$, the
sector~$H_2$ is $3$-dimensional, spanned by
$\{J^a_{-1}J^b_{-1}\mathbf{1}\}_{a,b \in \{+,-,0\}}$
modulo the Serre relations.  The pair-of-pants sewing kernel
restricted to this sector is a $3 \times 3$ matrix whose
entries are determined by the fusion rules
$V_1 \otimes V_1 \to V_1$ and the KZ connection at level~$1$.

The KZ equation at level~$1$ is
$\partial_z \Phi = \frac{\Omega}{(1+2)z}\Phi
= \frac{\Omega}{3z}\Phi$
(KZ normalisation with Sugawara denominator $k + h^\vee = 3$),
where $\Omega$ is the Casimir in $\mathfrak{sl}_2 \otimes
\mathfrak{sl}_2$.  The eigenvalues of the sewing kernel
$K_1$ on the $3$-dimensional sector are
$\{q, q^2, q^2\}$ (the level-$1$ WZW primary fields have
conformal weights $0$ and~$1/4$), so the trace-class condition
is satisfied with
$\|K_1|_{H_2}\|_1 = q + 2q^2$.

This is the gateway to a full interacting sewing
computation: the first case where the sewing kernel is a
non-diagonal matrix and the Fredholm expansion involves
genuine fusion data.
\end{example}

\begin{remark}[The conformal-block
 obstruction]\label{rem:conformal-block-obstruction}
For non-abelian chiral algebras at generic (non-integer) level,
the spaces of conformal blocks are infinite-dimensional; the
sewing kernel is an infinite-rank operator on these spaces.
HS-sewing guarantees that this operator is trace class, but the
explicit spectral decomposition requires detailed knowledge of
the conformal blocks (solutions of the KZ equation).  For
integer level, the conformal blocks are finite-dimensional
(Zhu's theorem), and the sewing kernel is a finite matrix at
each degree; this is the regime where explicit computation is
feasible.
\end{remark}

\subsection{Layer 2: metric independence of the ind-Hilbert
 factorisation}\label{subsec:gap-layer2}

Moriwaki's ind-Hilbert completion~\cite{Moriwaki26b} is defined
using a conformal metric on the curve; the Bergman inner product
on~$A^2(D)$ depends on the disk~$D$ and hence on the
conformal structure.  The algebraic bar complex, by contrast,
depends only on the OPE data and is independent of any metric.

\begin{definition}[Metric independence]
\label{def:metric-independence}
An ind-Hilbert factorisation is \emph{metrically independent}
if the isomorphism class of the completed sewing envelope
(as a topological vector space with continuous pair-of-pants
multiplication) is independent of the conformal metric used to
define the Bergman inner product.
\end{definition}

For the Heisenberg, metric independence is immediate: the Fock
space $\Sym A^2(D)$ depends on the disk~$D$, but any two disks
are conformally equivalent, and the isomorphism class of the
Bergman space is a conformal invariant.

For interacting algebras, metric independence is the assertion
that the analytic completion sees only the algebraic structure
(the OPE), not the auxiliary metric.  This is a functional-analytic
statement with no algebraic counterpart; it is the content of
the algebraic-to-analytic comparison theorem.

\begin{conjecture}[Metric independence]
\label{conj:metric-independence}
For every positive-energy chiral algebra satisfying HS-sewing,
the ind-Hilbert factorisation is metrically independent: the
sewing envelope~$\cA^{\mathrm{sew}}$ depends only on the
algebraic chiral core~$\cA_{\mathrm{alg}}$.
\end{conjecture}

\begin{remark}[Evidence for metric independence]
\label{rem:metric-evidence}
Three observations support the conjecture:
\begin{enumerate}[label=\textup{(\alph*)}]
\item For the Heisenberg, the Bergman space $A^2(D)$ depends
 on~$D$, but the abstract Hilbert space structure (separable,
 infinite-dimensional, with the canonical grading by polynomial
 degree) is a conformal invariant.  Any two conformal metrics
 give isomorphic Bergman spaces; the isomorphism is the
 conformal change-of-variables on holomorphic functions.
\item The algebraic bar differential depends only on OPE data
 (residues at collision diagonals); the OPE coefficients are
 independent of conformal metric.  If the bar differential
 is dense in the completed bar differential (as guaranteed by
 the analytic-algebraic comparison), the topology of the
 completion is determined by the algebraic data.
\item The sewing envelope is defined by
 seminorms~\eqref{eq:sewing-seminorm} that depend on bordism
 amplitudes.  Each amplitude is a composition of OPE operations
 (metric-independent) and sewing traces (metric-dependent only
 through the choice of conformal collar).  The collar dependence
 enters through the sewing parameter~$q$, which appears
 universally in all completions.
\end{enumerate}
The conjecture is that these observations extend to a proof for
all HS-sewing algebras: the sewing topology depends on~$q$ but
not on the conformal metric that determines the embedding of
the collar in the surface.
\end{remark}

\subsection{Layer 3: the coderived shadow at
 $g \geq 1$}\label{subsec:gap-layer3}

At genus $g = 0$ the bar complex satisfies $d^2 = 0$, and the
ordinary derived category is the correct receptacle.  At genus
$g \geq 1$, the fiberwise bar differential acquires curvature:
\begin{equation}\label{eq:curvature}
d_{\mathrm{fib}}^2
\;=\;
\kappa(\cA) \cdot \omega_g,
\end{equation}
where $\kappa(\cA)$ is the modular characteristic and
$\omega_g = c_1(\lambda)$ is the first Chern class of the
Hodge bundle on~$\overline\cM_g$.  For $\kappa \neq 0$
(which holds for every non-trivial chiral algebra in the
standard landscape), $d_{\mathrm{fib}}^2 \neq 0$ and
every object is acyclic in the ordinary derived category.

The correct receptacle is the coderived category~$\Dco$,
developed in the next section.  The passage from derived to
coderived is the categorified form of passage from the
equation~$d^2 = 0$ to the curved Maurer--Cartan equation
$d\alpha + \alpha^2 + \kappa\omega = 0$.

The analytic content of Layer~3 is: the sewing amplitudes that
converge by HS-sewing must be reinterpreted as evaluations of
coderived (not derived) functors.  The Fredholm determinant of
the Heisenberg is the prototype: the curvature
$d_{\mathrm{fib}}^2 = k \cdot \omega_1$ at genus~$1$ means that the
bar complex does not define an object in the derived category;
the Fredholm determinant
$\det(1 - T_q) = q^{-1/24}\eta(q)$ is the correct
coacyclic invariant.


\subsection{Interaction between the three
 layers}\label{subsec:gap-interaction}

The three layers are logically independent but interact in
concrete examples.

\begin{proposition}[Layer independence]
\label{prop:layer-independence}
\mbox{}
\begin{enumerate}[label=\textup{(\roman*)}]
\item Layer~1 (sewing envelope identification) does not imply
 Layer~2 (metric independence): an explicit Hilbert
 identification could depend on the choice of conformal metric.
\item Layer~2 (metric independence) does not imply Layer~3
 (coderived formulation): metric independence is a genus-$0$
 property of the completion, while the coderived category
 is forced by curvature at $g \geq 1$.
\item Layer~3 (coderived formulation) does not imply Layer~1:
 the coderived category provides the correct homotopy theory
 but does not identify the completion concretely.
\end{enumerate}
\end{proposition}

\begin{proof}
(i): the Heisenberg sewing envelope $\Sym A^2(D)$ depends on
the disk~$D$; the metric independence is a separate theorem
(conformal invariance of the Bergman space,
Proposition~\ref{prop:bergman-conformal}).

(ii): metric independence concerns the topology of the
completion; the coderived category concerns the homological
algebra of the curved differential.  The Heisenberg at
genus~$0$ is metrically independent but has $d^2 = 0$ (no
coderived issues); at genus~$1$ the coderived formulation
is needed but the metric independence is inherited from
genus~$0$.

(iii): the coderived category provides
$C^\bullet_{\mathrm{BV}} \simeq_{\Dco} B^{(g)}$ abstractly;
the identification of~$\cA^{\mathrm{sew}}$ with a specific
Hilbert space is additional data.
\end{proof}

For the standard landscape, the current status of each layer:

\begin{center}
\small
\renewcommand{\arraystretch}{1.15}
\begin{tabular}{lccc}
\toprule
Family & Layer 1 & Layer 2 & Layer 3 \\
\midrule
Heisenberg & proved & proved & proved \\
Lattice & proved & proved & proved \\
Affine KM & conjectural & conjectural & proved \\
Virasoro & conjectural & conjectural & proved \\
$\cW_N$ & conjectural & conjectural & proved \\
$\beta\gamma$ & conjectural & conjectural & conditional \\
\bottomrule
\end{tabular}
\end{center}

Layer~3 is proved for all standard families
(Theorem~\ref{thm:bv-bar-coderived}); Layers~1 and~2 are
proved only for the abelian families (Heisenberg and lattice).
The resolution of Layers~1 and~2 for non-abelian families is
the primary open problem of the analytic sewing programme.

\subsection{The sewing envelope as a nuclear
 space}\label{subsec:nuclear}

The sewing envelope carries a natural nuclear topology
(projective limit of Hilbert spaces~$H_q$ as $q \to 1^-$).
This nuclear structure provides additional functional-analytic
tools.

\begin{proposition}[Nuclear structure]
\label{prop:nuclear-structure}
Under HS-sewing, the sewing envelope $\cA^{\mathrm{sew}}$
is the projective limit
\[
\cA^{\mathrm{sew}}
\;=\;
\varprojlim_{q \to 1^-} H_q,
\]
and the transition maps $H_{q'} \to H_q$ (for $q' > q$) are
Hilbert--Schmidt.  In particular, $\cA^{\mathrm{sew}}$ is a
nuclear Fr\'echet space.
\end{proposition}

\begin{proof}
The weighted completions $H_q$ form a projective system: for
$q' > q$, the inclusion
$\cA_{\mathrm{alg}} \hookrightarrow H_q$ factors through
$H_{q'}$ because $q'^n > q^n$ for all $n \geq 1$.  The
transition map $H_{q'} \to H_q$ sends the basis vector
$e_{n,j} \in H_n$ to $(q/q')^n e_{n,j}$, with
Hilbert--Schmidt norm
$\sum_{n,j}(q/q')^{2n} = \sum_n \dim H_n (q/q')^{2n}$,
which is finite by subexponential sector growth (the ratio
$q/q' < 1$ provides exponential decay).  A projective limit
of Hilbert spaces with Hilbert--Schmidt transition maps is
nuclear~\cite{RS72}.
\end{proof}

The nuclear structure implies the Schwartz kernel theorem:
every continuous bilinear form on $\cA^{\mathrm{sew}}$ is
represented by a distribution kernel.  This provides the
correct framework for the pair-of-pants multiplication as a
continuous bilinear operation on the sewing envelope.


% ================================================================
% 8. THE CODERIVED CATEGORY AS CORRECT RECEPTACLE
% ================================================================

\section{The coderived category as correct receptacle}
\label{sec:coderived}

The genus $g \geq 1$ bar complex carries curvature
$d_{\mathrm{fib}}^2 = \kappa \cdot \omega_g$.  The ordinary
derived category is empty: every curved complex is acyclic.
The coderived category~$\Dco$ resolves this: curved
differentials are permitted, coacyclic complexes replace
acyclic ones, and the BV/bar comparison holds universally.

\subsection{Curved dg-modules}\label{subsec:curved-dg}

\begin{definition}[Curved dg-module]\label{def:curved-dg}
A \emph{curved dg-module} over a graded algebra~$R$ is a
graded $R$-module~$M$ equipped with a degree-$1$
endomorphism~$d$ satisfying
\[
d^2 \;=\; \mu \cdot \id_M
\]
for some central element~$\mu \in R$ (the \emph{curvature}).
When $\mu = 0$ this reduces to an ordinary dg-module.
\end{definition}

The genus-$g$ bar complex $B^{(g)}(\cA)$ is a curved
dg-module over
$R = H^*(\overline\cM_g)$ with curvature
$\mu = \kappa(\cA) \cdot \omega_g$.

\begin{warning}[Acyclicity in the ordinary sense]
\label{warn:acyclicity}
If $d^2 = \mu \neq 0$, then for every
cocycle~$x$ ($dx = 0$) we have
$\mu \cdot x = d^2 x = d(dx) = 0$; if $\mu$ is
not a zero divisor, this forces $x = 0$.  In particular,
$H^*(M,d) = 0$ whenever $\mu$ is a non-zero-divisor.  The
ordinary derived category of curved dg-modules is therefore
trivial when $\kappa \neq 0$.
\end{warning}

\subsection{The coderived category}\label{subsec:coderived-def}

The coderived category replaces acyclicity with coacyclicity.

\begin{definition}[Coacyclic complex]
\label{def:coacyclic}
A curved dg-module~$M$ is \emph{coacyclic} if it lies in the
smallest thick subcategory of the homotopy category of curved
dg-modules containing all totalizations of short exact
sequences.  Equivalently, $M$ is coacyclic if every morphism
from~$M$ to an injective curved dg-module is null-homotopic.
\end{definition}

\begin{definition}[Coderived category]
\label{def:coderived-category}
The \emph{coderived category}~$\Dco(R\text{-}\mathrm{mod}_\mu)$
is the Verdier quotient of the homotopy category of curved
dg-modules by the thick subcategory of coacyclic complexes.
\end{definition}

The coderived category is the correct homotopy category for
curved complexes: it retains the meaningful homological
information that the ordinary derived category discards.

\begin{remark}[Positselski's theory]
\label{rem:positselski}
The coderived category was introduced by
Positselski~\cite{Positselski11} in the context of curved
dg-algebras and matrix factorisations.  The totalisation of a
short exact sequence of curved complexes need not be acyclic
(because $d^2 \neq 0$), but it is always coacyclic:
coacyclicity is the correct generalisation of acyclicity to the
curved setting.
\end{remark}

\subsection{The curvature bicomplex}\label{subsec:curvature-bicomplex}

The curved bar complex $B^{(g)}(\cA)$ carries a natural
bicomplex structure.  Write the total differential as
\[
D_g \;=\; d_{\mathrm{fib}} + \nabla^{\mathrm{GM}}_g,
\]
where $d_{\mathrm{fib}}$ is the fiberwise bar differential
(with curvature $d_{\mathrm{fib}}^2 = \kappa \cdot \omega_g$)
and $\nabla^{\mathrm{GM}}_g$ is the Gauss--Manin connection.
The total differential satisfies
\begin{equation}\label{eq:total-flat}
D_g^2
\;=\;
d_{\mathrm{fib}}^2
+ [d_{\mathrm{fib}}, \nabla^{\mathrm{GM}}_g]
+ (\nabla^{\mathrm{GM}}_g)^2
\;=\;
\kappa\omega_g - \kappa\omega_g
\;=\; 0:
\end{equation}
flatness is restored by coupling the bar differential to the
Hodge bundle.  The curvature bicomplex is the bigraded object
with horizontal differential~$d_{\mathrm{fib}}$ and vertical
differential~$\nabla^{\mathrm{GM}}_g$.

\begin{proposition}[Totalisation in $\Dco$]
\label{prop:totalisation}
The totalisation of the curvature bicomplex of~$B^{(g)}(\cA)$
is an object of the coderived category~$\Dco$.  When~$\cA$
satisfies HS-sewing, the associated partition function is the
trace of the totalised sewing operator and coincides with the
absolutely convergent Fredholm expansion of
Corollary~\ref{cor:partition-convergence}.
\end{proposition}

\begin{proof}
The totalisation of the bicomplex
$(d_{\mathrm{fib}}, \nabla^{\mathrm{GM}}_g)$ produces a
complex with $D_g^2 = 0$ by~\eqref{eq:total-flat}, hence an
honest object of the derived category of the total complex.
The passage to the coderived category of the curved
component~$d_{\mathrm{fib}}$ (with
$d_{\mathrm{fib}}^2 = \kappa\omega_g$) is by the standard
Positselski comparison: totalisation of the bicomplex
represents the curved component as an object of~$\Dco$
because the vertical differential~$\nabla^{\mathrm{GM}}_g$
absorbs the curvature.

Under HS-sewing, the trace-class property of
Proposition~\ref{prop:closed-trace-class} ensures that the
trace of the totalised sewing operator is an absolutely
convergent number; this trace coincides with the partition
function computed from the Fredholm expansion, because both
are the same composition of pair-of-pants operators evaluated
on the completed Hilbert space~$H_q$.
\end{proof}

\subsection{The BV/bar comparison in the coderived
 category}\label{subsec:bv-bar}

The BV complex of a chiral algebra on a Riemann
surface~$\Sigma_g$ is the Costello--Gwilliam perturbative
quantisation~\cite{CG17}; the bar complex $B^{(g)}(\cA)$ is
the algebraic chain model.  At genus~$0$, the two are
chain-level quasi-isomorphic.  At genus~$g \geq 1$, curvature
intervenes: the BV propagator acquires a harmonic component
(the Arakelov Green function) that produces corrections
invisible to the algebraic bar differential.

\begin{theorem}[BV${}={}$bar in the coderived category]
\label{thm:bv-bar-coderived}
Let~$\cA$ be a chirally Koszul algebra.  For each genus
$g \geq 1$, let
\[
f_g \colon
C^\bullet_{\mathrm{BV}}(\cA, \Sigma_g)
\longrightarrow
B^{(g)}(\cA)
\]
be the comparison morphism obtained by replacing the BV
propagator with its Hodge decomposition relative to the bar
propagator.
\begin{enumerate}[label=\textup{(\roman*)}]
\item For every genus $g \geq 0$,
\[
C^\bullet_{\mathrm{BV}}(\cA, \Sigma_g)
\;\simeq_{\Dco}\;
B^{(g)}(\cA).
\]
\item If $\cA$ lies in class~$\mathsf{G}$ or~$\mathsf{L}$,
 then~$f_g$ is a weak equivalence of filtered curved models.
\item If $\cA$ lies in class~$\mathsf{C}$ and harmonic
 decoupling holds, then~$f_g$ is a weak equivalence.
\item If $\cA$ lies in class~$\mathsf{M}$, then for every
 $r \geq 4$ the harmonic discrepancy satisfies
 \[
 \delta_r^{\mathrm{harm}}
 \;=\;
 c_r(\cA) \cdot m_0^{\lfloor r/2 \rfloor - 1}
 \]
 and the cone of~$f_g$ is coacyclic.
\end{enumerate}
\end{theorem}

\begin{proof}[Proof sketch]
At genus~$0$, the comparison is the chain-level
quasi-isomorphism: $d_{\mathrm{fib}}^2 = 0$ and the harmonic
component of the propagator vanishes on~$\bC$.

For $g \geq 1$, Hodge decomposition writes the BV propagator
as the sum of the bar propagator (holomorphic, $d\log E(z,w)$,
weight~$1$) and the harmonic correction
$h_g = (2\pi)^{-1}(\operatorname{Im}\Omega)^{-1}_{ab}\,
d\bar z_a \wedge dz_b$.  The harmonic correction produces a
discrepancy between the BV and bar differentials.

For classes~G and~L (Heisenberg, affine Kac--Moody): the
harmonic correction to the cubic vertex vanishes by Jacobi
(for~L) or by absence of interaction vertices (for~G).  The
comparison is a weak equivalence of filtered curved models.

For class~C ($\beta\gamma$): composite free-field
factorisation reduces the harmonic correction to a product of
Heisenberg corrections; each factor vanishes by the G-class
argument, conditional on harmonic decoupling.

For class~M (Virasoro, $\cW_N$): the quartic harmonic
discrepancy
$\delta_4^{\mathrm{harm}} \propto
Q^{\mathrm{contact}} \cdot \kappa/\operatorname{Im}(\tau)$
is not cancelled by the Fay trisecant identity.  The naive
chain-level comparison fails.  The full harmonic discrepancy
is a sum of curvature powers
$\delta_r^{\mathrm{harm}}
= c_r(\cA) \cdot m_0^{\lfloor r/2\rfloor - 1}$,
and the comparison cone is coacyclic in the sense of
Definition~\ref{def:coacyclic}: every morphism from the cone
to an injective curved module is null-homotopic, because the
discrepancy is divisible by the curvature element~$m_0$.

The coderived equivalence~(i) follows: the cone of~$f_g$ is
coacyclic for every class, and a morphism with coacyclic cone
is an isomorphism in~$\Dco$.
\end{proof}

\begin{remark}[Class-by-class obstruction profile]
\label{rem:class-by-class}
The four shadow classes of the standard landscape exhibit a
strict hierarchy of analytic regularity:
\begin{center}
\small
\renewcommand{\arraystretch}{1.15}
\begin{tabular}{lcl}
\toprule
Class & Chain-level status & Mechanism \\
\midrule
$\mathsf{G}$ (Heisenberg) & weak equivalence
 & no interaction vertices \\
$\mathsf{L}$ (affine KM) & weak equivalence
 & Jacobi kills cubic harmonic term \\
$\mathsf{C}$ ($\beta\gamma$) & conditional
 & composite free-field factorisation \\
$\mathsf{M}$ (Virasoro, $\cW_N$) & coderived only
 & harmonic correction curvature-divisible \\
\bottomrule
\end{tabular}
\end{center}
For classes~G and~L the chain-level and coderived comparisons
coincide.  For class~M the chain-level comparison genuinely
fails, and the coderived category is essential.
\end{remark}

\subsection{The harmonic discrepancy mechanism}
\label{subsec:harmonic-mechanism}

The key to the coderived comparison is the factorisation of
the harmonic discrepancy through curvature powers.  The
mechanism is the Hodge decomposition of the Bergman kernel.

On a genus-$g$ surface~$\Sigma_g$ with period
matrix~$\Omega$, the Bergman kernel (the bar propagator) is
\begin{equation}\label{eq:bergman-g}
B(z,w)
\;=\;
d_z d_w \log E(z,w)
\;+\;
2\pi i\, (\operatorname{Im}\Omega)^{-1}_{ab}\,
\omega_a(z)\,\omega_b(w),
\end{equation}
where $E(z,w)$ is the prime form and
$\omega_1,\ldots,\omega_g$ is a normalised basis of abelian
differentials.  The second term is the harmonic correction.

The BV propagator~$P_g$ differs from the bar
propagator~$d\log E(z,w)$ by this harmonic term:
\[
P_g(z,w)
\;=\;
d\log E(z,w)
\;+\;
h_g(z,w),
\qquad
h_g
\;=\;
\frac{1}{2\pi}\,
(\operatorname{Im}\Omega)^{-1}_{ab}\,
d\bar z_a \wedge dz_b.
\]
Each insertion of~$h_g$ into a Feynman graph vertex produces
a harmonic correction to the BV amplitude.  At the quartic
level ($r = 4$), the correction is
\begin{equation}\label{eq:quartic-obstruction}
\delta_4^{\mathrm{harm}}
\;=\;
Q^{\mathrm{contact}}(\cA)
\cdot \frac{\kappa(\cA)}{\operatorname{Im}(\tau)},
\end{equation}
where $Q^{\mathrm{contact}}$ is the quartic contact invariant
of~$\cA$.  For class~G and~L, $Q^{\mathrm{contact}} = 0$.
For class~M, $Q^{\mathrm{contact}} \neq 0$ and
$\delta_4^{\mathrm{harm}}$ is a genuine obstruction to the
chain-level comparison.

The crucial observation is that every~$\delta_r^{\mathrm{harm}}$
is divisible by~$m_0 = \kappa \cdot \omega_g$ (the curvature
element of the curved bar differential).  This divisibility
is the mechanism by which the comparison cone becomes
coacyclic: in the coderived category, elements divisible by the
curvature are trivial.

\subsection{Coderived versus derived: a comparison}
\label{subsec:coderived-comparison}

The passage from derived to coderived captures a precise
mathematical phenomenon.

\begin{center}
\small
\renewcommand{\arraystretch}{1.15}
\begin{tabular}{lcc}
\toprule
 & \textbf{Derived} $D(R)$ & \textbf{Coderived}
  $\Dco(R)$ \\
\midrule
Curvature & $d^2 = 0$ required & $d^2 = \mu$ permitted \\
Acyclicity notion & $H^*(M) = 0$ & coacyclic \\
Genus-$0$ bar & object of $D(R)$ & object of $\Dco(R)$ \\
Genus-$g$ bar ($\kappa \neq 0$) & empty & well-defined \\
BV/bar comparison & chain-level (G,L) & all classes \\
Partition function & formal series & convergent number \\
\bottomrule
\end{tabular}
\end{center}

The coderived category is the categorified avatar of the
curved Maurer--Cartan equation.  The passage from
$d^2 = 0$ to $d\alpha + \alpha^2 + \kappa\omega = 0$
is the passage from derived to coderived: the curvature
becomes data rather than obstruction.

\subsection{Genus-1 illustration}\label{subsec:genus-1-coderived}

The genus-$1$ case illustrates the three categories concretely.

At genus~$1$, the bar complex $B^{(1)}(\cA)$ has
curvature $d_{\mathrm{fib}}^2 = \kappa(\cA) \cdot \omega_1$,
where $\omega_1 = c_1(\lambda)$ is the first Chern class
of the Hodge bundle on~$\overline\cM_{1,1}$.  The total
differential $D_1 = d_{\mathrm{fib}} + \nabla^{\mathrm{GM}}$
satisfies $D_1^2 = 0$.

\begin{example}[Heisenberg at genus~$1$]
\label{ex:heisenberg-genus1-coderived}
For $\cH_k$ with $\kappa = k$:
\begin{enumerate}[label=\textup{(\alph*)}]
\item \emph{Derived category:} $d_{\mathrm{fib}}^2
 = k \cdot \omega_1 \neq 0$ for $k \neq 0$.
 Every object in the curved complex is acyclic:
 $H^*(B^{(1)}(\cH_k), d_{\mathrm{fib}}) = 0$.
\item \emph{Coderived category:} the totalisation of the
 bicomplex $(d_{\mathrm{fib}}, \nabla^{\mathrm{GM}})$ is a
 well-defined object of~$\Dco$.  The trace of the sewing
 operator gives
 $Z_1 = (\operatorname{Im}\tau)^{-k/2}\,|\eta(\tau)|^{-2k}
 \neq 0$.
\item \emph{Total category:} $D_1^2 = 0$; the total complex
 $B^{(1)}(\cH_k)$ is an honest dg-module over
 $H^*(\overline\cM_{1,1})$ and lives in the ordinary derived
 category of the total complex.
\end{enumerate}
The partition function~$Z_1$ is a coderived invariant
(well-defined in~$\Dco$ of the curved fiberwise complex) and
simultaneously a derived invariant (well-defined in~$D$ of the
total complex).  The two perspectives are related by the
totalisation functor of
Proposition~\ref{prop:totalisation}.
\end{example}

\begin{example}[Virasoro at genus~$1$]
\label{ex:virasoro-genus1-coderived}
For $\Vir_c$ with $\kappa = c/2$:
\begin{enumerate}[label=\textup{(\alph*)}]
\item The curvature is
 $d_{\mathrm{fib}}^2 = (c/2)\cdot\omega_1$.
\item The partition function is
 $Z_1(\Vir_c;\,\tau)
 = (\operatorname{Im}\tau)^{-c/4}\,
 |\eta(\tau)|^{-2(c-1)}$
 (the vacuum character of the Virasoro algebra).
\item The BV/bar comparison: the harmonic discrepancy at the
 quartic level is
 $\delta_4^{\mathrm{harm}}
 = Q^{\mathrm{contact}} \cdot (c/2)/\operatorname{Im}(\tau)$.
 For $c \neq 0$ (non-trivial Virasoro), this is nonzero; the
 chain-level BV/bar comparison fails.  The cone is coacyclic
 by Theorem~\ref{thm:bv-bar-coderived}(iv), so the coderived
 comparison holds.
\end{enumerate}
The Virasoro at genus~$1$ is the simplest class-M example: the
harmonic discrepancy is explicit, the coacyclicity is verifiable,
and the coderived category is genuinely necessary.
\end{example}

\subsection{Convergent versus formal partition
 functions}\label{subsec:convergent-formal}

The three levels of structure (algebraic, coderived, analytic)
produce three notions of partition function:

\begin{enumerate}[label=\textup{(\roman*)}]
\item \emph{Formal:} the genus-$g$ component of the MC element
 $\Theta_\cA$, evaluated on the fundamental class
 $[\overline\cM_g]$.  This is a formal power series in the
 sewing parameters $q_1,\ldots,q_g$.  It exists for every
 chiral algebra with $D^2 = 0$.
\item \emph{Coderived:} the trace of the totalised sewing
 operator in the coderived category.  This is a formal power
 series with the additional property that each coefficient is a
 coderived invariant (independent of the choice of resolution
 up to coacyclic error).  It exists for every positive-energy
 chiral algebra.
\item \emph{Analytic:} an absolutely convergent complex number,
 the trace of a trace-class operator on the weighted Hilbert
 completion~$H_q$.  This exists under HS-sewing.
\end{enumerate}

\begin{proposition}[Compatibility of the three notions]
\label{prop:three-partition-functions}
For a chirally Koszul algebra satisfying HS-sewing:
\begin{enumerate}[label=\textup{(\alph*)}]
\item the formal partition function is the Taylor expansion of
 the analytic partition function around the cusp
 $q_j \to 0$;
\item the coderived partition function is the image of the
 analytic partition function under the forgetful functor from
 the Hilbert-completed category to the coderived category;
\item all three notions agree on the domain of convergence.
\end{enumerate}
\end{proposition}

\begin{proof}
(a): the formal MC element $\Theta_\cA$ is the generating
function for the graph-sum amplitudes, which are the Taylor
coefficients of the partition function at the cusp.  Under
HS-sewing, the Taylor series converges absolutely for
$|q_j| < 1$ (Corollary~\ref{cor:partition-convergence}), so
the formal and analytic partition functions agree on the
domain of convergence.

(b): the forgetful functor sends a trace-class operator on~$H_q$
to its underlying curved dg-module structure in~$\Dco$; the
trace (a convergent number) is sent to the same number viewed
as a coderived invariant.

(c): on the domain of convergence, all three are the same
complex number, computed by the same graph-sum with the same
absolutely convergent Fredholm expansion.
\end{proof}

\subsection{The analytic Koszul pair}\label{subsec:analytic-koszul}

The algebraic Koszul pair $(\cA, \cA^!)$ controls Theorems~A
through~H of the modular Koszul duality programme
at the algebraic level.  The analytic extension is the
following.

\begin{definition}[Analytic Koszul pair]
\label{def:analytic-koszul-pair}
An \emph{analytic Koszul pair} is a pair
$(\cA, \cA^!)^{\mathrm{an}}$ of chirally Koszul algebras,
both satisfying HS-sewing, for which the analytic bar-cobar
adjunction restricts to a Quillen equivalence on
sewing-complete categories:
\[
\barB^{\mathrm{an}} \colon
\cA^{\mathrm{sew}}\text{-}\mathrm{mod}
\rightleftarrows
(\cA^!)^{\mathrm{sew}}\text{-}\mathrm{comod}
\;\colon\!
\Omega^{\mathrm{an}}.
\]
\end{definition}

Every algebraic Koszul pair in the standard landscape is
expected to extend to an analytic Koszul pair; the HS-sewing
theorem provides the analytic input for both~$\cA$
and~$\cA^!$.

\begin{conjecture}[Analytic realisation criterion]
\label{conj:analytic-realisation}
HS-sewing on~$\cA$ plus dual HS-sewing on~$\cA^!$ suffices
for the analytic Koszul pair to exist: the analytic bar-cobar
adjunction restricts to a Quillen equivalence on
sewing-complete categories.
\end{conjecture}

The verification of this conjecture for the Heisenberg
($\cH_k^! = \Sym^{\mathrm{ch}}(V^*) \neq \cH_{-k}$) and
for lattice vertex algebras is the natural test case.

\begin{remark}[Koszul complementarity and the analytic
 pair]\label{rem:koszul-complementarity-analytic}
The Koszul complementarity theorem (Theorem~C of the modular
Koszul duality programme) asserts that
$\kappa(\cA) + \kappa(\cA^!) = K(\cA)$
(the Koszul conductor): $K = 0$ for KM/Heis/lattice/free;
$K = 13$ for Virasoro; $K = 250/3$ for $\cW_3$.  The
analytic content is: both $\cA$ and $\cA^!$ satisfy HS-sewing
(because both are in the standard landscape), and the
complementarity relation constrains the sewing data:
\begin{equation}\label{eq:complementarity-sewing}
S_\cA(u) + S_{\cA^!}(u)
\;=\;
K \cdot \zeta(u)\,\zeta(u+1)
\;+\;
\text{lower-order terms},
\end{equation}
where the leading term involves the Koszul conductor.  For
the Virasoro pair $(\Vir_c, \Vir_{26-c})$:
$S_{\Vir_c}(u) + S_{\Vir_{26-c}}(u)
= 2\zeta(u+1)(\zeta(u) - 1)$,
independent of~$c$; the dependence on~$c$ cancels in the sum,
reflecting the topological nature of the Koszul conductor.
\end{remark}

\begin{remark}[The analytic bar coalgebra]
\label{rem:analytic-bar-coalgebra}
The algebraic bar coalgebra
$B(\cA) = T^c(s^{-1}\bar\cA)$ carries the deconcatenation
coproduct $\Delta \colon B(\cA) \to B(\cA) \otimes B(\cA)$.
Under HS-sewing, the bar coalgebra completes to the
\emph{analytic bar coalgebra}
$B^{\mathrm{an}}(\cA)$, defined as the graph-norm closure
of $B(\cA)$ in the Hilbert-space topology induced by~$H_q$.
The coproduct extends continuously:
$\Delta^{\mathrm{an}} \colon B^{\mathrm{an}}(\cA) \to
B^{\mathrm{an}}(\cA)
\widehat{\otimes} B^{\mathrm{an}}(\cA)$,
where $\widehat{\otimes}$ is the completed Hilbert-space
tensor product.

The analytic bar coalgebra is the correct object for the
analytic bar-cobar adjunction: it carries both the algebraic
coalgebra structure (deconcatenation) and the analytic
topology (trace-class sewing), and the cobar functor
$\Omega^{\mathrm{an}}$ applied to
$B^{\mathrm{an}}(\cA)$ recovers the sewing
envelope~$\cA^{\mathrm{sew}}$.
\end{remark}

\subsection{The curvature--sewing
 interplay}\label{subsec:curvature-sewing}

The coderived category and the HS-sewing condition interact
at genus $g \geq 1$ through the curvature.

\begin{proposition}[Curvature and trace class]
\label{prop:curvature-trace}
Let $\cA$ satisfy HS-sewing and let $\kappa = \kappa(\cA)$.
At genus $g \geq 1$, the curvature element
$m_0 = \kappa \cdot \omega_g$ acts on the weighted
completion~$H_q$ as a bounded operator of norm
$\|m_0\|_{\mathrm{op}} = |\kappa| \cdot \|\omega_g\|$,
where $\|\omega_g\|$ is the sup-norm of the Hodge class
representative on the moduli space.  The iterated curvatures
$m_0^k$ remain bounded for all~$k$, and the harmonic
discrepancy $\delta_r^{\mathrm{harm}}
= c_r \cdot m_0^{\lfloor r/2\rfloor - 1}$
is a bounded operator on~$H_q$ for every~$r$.
\end{proposition}

\begin{proof}
The curvature element $m_0 = \kappa \cdot \omega_g$ is a
scalar operator (multiplication by a cohomology class)
tensored with the identity on the state space.  As such, its
operator norm on~$H_q$ is $|\kappa| \cdot \|\omega_g\|$,
independent of~$q$.  Powers are similarly bounded:
$\|m_0^k\| = |\kappa|^k \cdot \|\omega_g\|^k$.  The
harmonic discrepancy $\delta_r^{\mathrm{harm}}$ is a
finite sum of such terms with coefficients $c_r(\cA)$
determined by the OPE data; each term is bounded.
\end{proof}

The proposition shows that the curvature does not disrupt
the trace-class property of the sewing operators: the
HS-sewing condition and the curved differential coexist in
the completed Hilbert space.  The coderived category
provides the correct homological framework; the HS-sewing
condition provides the correct functional-analytic framework;
together they deliver convergent genus-$g$ partition functions
for curved bar complexes.

\begin{proposition}[Genus-$g$ partition function as coderived
 trace]\label{prop:coderived-trace}
Under HS-sewing, the genus-$g$ partition function is
simultaneously:
\begin{enumerate}[label=\textup{(\alph*)}]
\item the trace of a trace-class operator on~$H_q$
 (analytic);
\item an evaluation of the MC element
 $\Theta_\cA$ on $[\overline\cM_g]$ (algebraic);
\item a coderived invariant of the curved bar
 complex~$B^{(g)}(\cA)$ in~$\Dco$ (homological).
\end{enumerate}
All three produce the same complex number on the domain of
convergence.
\end{proposition}

\begin{proof}
(a) is Corollary~\ref{cor:partition-convergence}.
(b) is the definition of the MC evaluation (the graph-sum
expansion).  (c) is
Proposition~\ref{prop:three-partition-functions}.  The
equality of the three numbers follows from the compatibility
statement of
Proposition~\ref{prop:three-partition-functions}(c).
\end{proof}


% ================================================================
% CONCLUSION
% ================================================================

\section*{Conclusion}

The results of this paper organise as follows.

The sewing envelope (Definition~\ref{def:sewing-envelope})
provides the universal locally convex completion for convergent
amplitudes.  The HS-sewing condition
(Definition~\ref{def:hs-sewing}) is a single Hilbert--Schmidt
estimate that implies trace-class closed amplitudes at all
genera.  For the Heisenberg, the sewing programme is
completely explicit: Bergman space, second quantisation,
Fredholm determinant
(Theorems~\ref{thm:heisenberg-sewing}
and~\ref{thm:heisenberg-one-particle}).  For lattice vertex
algebras, the programme extends by combining the Heisenberg
Fredholm determinant with convergent theta functions
(Theorem~\ref{thm:lattice-sewing}).  The general HS-sewing
theorem (Theorem~\ref{thm:general-hs-sewing}) verifies the
criterion for the entire standard landscape.

Three layers separate the convergence statement from a complete
analytic theory: the sewing envelope for interacting algebras
(explicit Hilbert identification open), metric independence of
the ind-Hilbert factorisation (conjectural), and the coderived
shadow at $g \geq 1$ (the BV/bar comparison in~$\Dco$ is
proved for all four shadow classes,
Theorem~\ref{thm:bv-bar-coderived}).

The coderived category is the correct receptacle for the curved
bar complex: curvature is data, not obstruction; the harmonic
discrepancy for class~M is coacyclic (divisible by the
curvature element); and the partition functions computed by
HS-sewing are coderived invariants.

Five problems define the frontier.

\paragraph{The sewing envelope of $V_1(\mathfrak{sl}_2)$.}
The simplest interacting case.  The conformal blocks are
finite-dimensional at each degree; the sewing kernel is a
finite matrix computable from the KZ connection at level~$1$.
The goal is an explicit Fredholm expansion involving Weyl-orbit
data, providing the template for all non-abelian sewing
envelopes.

\paragraph{Fredholm determinants beyond Heisenberg.}
For non-abelian Kac--Moody and Virasoro, the closed amplitudes
are trace class but are not yet known to be Fredholm
determinants of a single one-particle operator.  The
Heisenberg identification suggests the existence of an
analogous second-quantisation structure for every chiral
algebra with quadratic OPE; finding the correct one-particle
Hilbert space is the natural next step.

\paragraph{The analytic bar-cobar adjunction.}
The algebraic bar-cobar adjunction
$\barB \dashv \Omega$ is the structural spine of the
modular Koszul duality programme.  Conjecture~\ref{conj:analytic-realisation}
predicts that HS-sewing on both sides
($\cA$ and $\cA^!$) suffices for the analytic adjunction.
Verification for the Heisenberg pair
$(\cH_k, \Sym^{\mathrm{ch}}(V^*))$ and the Virasoro pair
$(\Vir_c, \Vir_{26-c})$ are the natural test cases.  For
Virasoro, the Koszul dual has central charge $c' = 26 - c$ and
$\kappa(\Vir_c) + \kappa(\Vir_{26-c}) = c/2 + (26-c)/2 = 13$;
both sides satisfy HS-sewing.

\paragraph{Metric independence.}
Conjecture~\ref{conj:metric-independence} predicts that the
sewing envelope is independent of the conformal metric.
A proof for the Heisenberg (where the statement is elementary)
and for lattice vertex algebras (where it follows from the
conformal invariance of the theta function) would establish the
pattern; the general case requires new functional-analytic
techniques.

\paragraph{Positivity and modular invariance.}
Corollary~\ref{cor:hs-positivity} establishes a positivity
cone for HS-sewing amplitudes; it does not directly imply
$\mathrm{SL}_2(\bZ)$ covariance of genus-$1$ characters.
Combining the analytic sewing of this paper with Zhu's modular
invariance theorem~\cite{Zhu96} and Huang's higher-genus
extension~\cite{Huang05} is the natural route to a unified
analytic modular programme.

\bigskip

The algebraic bar complex and the analytic sewing envelope are
two facets of the same object: the bar complex is the formal
skeleton; the sewing envelope is the convergent flesh.  The
coderived category mediates between them, absorbing the
curvature that obstructs the ordinary derived category at
$g \geq 1$.  The modular Koszul duality programme operates at
the algebraic level; the analytic sewing programme developed
here is the bridge to computable, convergent partition functions.


% ================================================================
% APPENDIX: FUNCTIONAL ANALYSIS BACKGROUND
% ================================================================

\appendix

\section{Functional analysis background}
\label{app:functional-analysis}

We collect the functional-analytic facts used throughout.

\subsection{Trace-class and Hilbert--Schmidt
 operators}\label{app:trace-hs}

Let $H$ be a separable Hilbert space.

\begin{definition}\label{def:trace-class}
A bounded operator $T \colon H \to H$ is:
\begin{enumerate}[label=\textup{(\alph*)}]
\item \emph{Hilbert--Schmidt} if
 $\|T\|_{\HS}^2 := \sum_j \|Te_j\|^2 < \infty$
 for some (hence every) orthonormal basis $\{e_j\}$;
\item \emph{trace class} if
 $\|T\|_1 := \sum_j \langle e_j, |T| e_j\rangle < \infty$
 where $|T| = (T^*T)^{1/2}$;
\item \emph{nuclear} if it admits a representation
 $T = \sum_j \lambda_j \langle f_j, \cdot\rangle g_j$
 with $\sum_j |\lambda_j| < \infty$.
\end{enumerate}
The three conditions (a), (b), (c) are equivalent for
positive operators; in general,
trace class $\subset$ Hilbert--Schmidt $\subset$ compact.
\end{definition}

\begin{proposition}[Composition and trace]
\label{prop:hs-composition}
\mbox{}
\begin{enumerate}[label=\textup{(\roman*)}]
\item The composition of two Hilbert--Schmidt operators is
 trace class.
\item If $T$ is trace class, then
 $\Tr(T) = \sum_j \langle e_j, Te_j\rangle$ converges
 absolutely and is independent of the basis.
\item If $T$ is Hilbert--Schmidt and $S$ is bounded, then
 $ST$ and $TS$ are Hilbert--Schmidt.
\end{enumerate}
\end{proposition}

These standard facts (see Reed--Simon~\cite{RS72}) underpin
the trace-class theorem
(Proposition~\ref{prop:closed-trace-class}): each pair of
pants contributes a Hilbert--Schmidt operator; each sewing
circle composes two Hilbert--Schmidt operators into a
trace-class operator.

\subsection{Fredholm determinants}\label{app:fredholm}

\begin{definition}\label{def:fredholm-det}
For a trace-class operator~$K$ on a separable Hilbert space,
the \emph{Fredholm determinant} is
\[
\det(1 - K)
\;=\;
\prod_{j \geq 1}(1 - \lambda_j),
\]
where $\{\lambda_j\}$ are the eigenvalues of~$K$ counted with
algebraic multiplicity.  The product converges absolutely
whenever~$K$ is trace class.
\end{definition}

\begin{proposition}[Properties of the Fredholm determinant]
\label{prop:fredholm-properties}
Let $K$ be trace class.
\begin{enumerate}[label=\textup{(\roman*)}]
\item $|\det(1 - K)| \leq \exp(\|K\|_1)$.
\item $\log\det(1 - K)
 = -\sum_{m \geq 1}\frac{1}{m}\Tr(K^m)$
 whenever $\|K\| < 1$.
\item $\det(1 - K) \neq 0$ if and only if
 $1 \notin \mathrm{spec}(K)$.
\item The map $K \mapsto \det(1 - K)$ is continuous in the
 trace norm.
\end{enumerate}
\end{proposition}

For the Heisenberg sewing, the operator~$K_g$ has
eigenvalues~$q^n$ ($n \geq 1$) at genus~$1$; the Fredholm
determinant is $\det(1 - K_1) = \prod_{n \geq 1}(1-q^n)
= q^{-1/24}\eta(q)$.

\subsection{The Bergman space}\label{app:bergman}

\begin{definition}\label{def:bergman-space}
The \emph{Bergman space} $A^2(\Omega)$ of a bounded domain
$\Omega \subset \bC$ is the Hilbert space of
square-integrable holomorphic functions:
\[
A^2(\Omega)
\;=\;
\bigl\{f \in \cO(\Omega) :
\|f\|^2 := \textstyle\int_\Omega |f|^2\,dA < \infty
\bigr\}.
\]
The Bergman kernel
$K_\Omega(z,w) = \sum_j e_j(z)\overline{e_j(w)}$
(with $\{e_j\}$ any orthonormal basis) is the reproducing
kernel: $f(z) = \int_\Omega K_\Omega(z,w) f(w)\,dA(w)$.
\end{definition}

For the unit disk,
$K_D(z,w) = \frac{1}{\pi}(1-z\bar w)^{-2}$
and the orthonormal basis is
$e_n(z) = \sqrt{(n+1)/\pi}\,z^n$
for $n \geq 0$.  The Bergman space of the disk is the
one-particle Hilbert space for the Heisenberg sewing
(Definition~\ref{def:theta-map}).

\begin{proposition}[Conformal invariance]
\label{prop:bergman-conformal}
If $\varphi \colon \Omega_1 \to \Omega_2$ is a biholomorphism,
then the map $f \mapsto (f \circ \varphi)\cdot\varphi'$ is a
unitary isomorphism
$A^2(\Omega_2) \xrightarrow{\sim} A^2(\Omega_1)$.  In
particular, the Bergman space of any simply connected domain is
unitarily isomorphic to~$A^2(D)$.
\end{proposition}

This conformal invariance is the mechanism behind the metric
independence of the Heisenberg sewing envelope: any choice of
conformal collar gives an isomorphic Bergman space.

\subsection{Siegel modular forms}\label{app:siegel}

The Siegel upper half-space of degree~$g$ is
\[
\fH_g
\;=\;
\{\Omega \in \mathrm{Mat}_{g \times g}(\bC) :
\Omega^T = \Omega,\;
\operatorname{Im}\Omega > 0\}.
\]
The symplectic group
$\mathrm{Sp}_{2g}(\bZ)$ acts on~$\fH_g$ by
$\begin{pmatrix}A&B\\C&D\end{pmatrix}
\cdot \Omega = (A\Omega+B)(C\Omega+D)^{-1}$.

\begin{definition}\label{def:siegel-theta}
The \emph{Siegel theta function} of an even lattice~$\Lambda$
of rank~$r$ is
\[
\Theta_\Lambda(\Omega)
\;=\;
\sum_{\lambda \in \Lambda^g}
e^{i\pi\lambda^T\Omega\lambda},
\qquad \Omega \in \fH_g.
\]
For unimodular~$\Lambda$, $\Theta_\Lambda$ is a Siegel modular
form of weight~$r/2$ with respect to
$\mathrm{Sp}_{2g}(\bZ)$.
\end{definition}

At genus~$1$, $\fH_1 = \fH$ is the classical upper half-plane,
$\Omega = \tau$, and the Siegel theta function reduces to the
Jacobi theta function
$\Theta_\Lambda(\tau)
= \sum_{\lambda \in \Lambda} q^{\langle\lambda,\lambda\rangle/2}$
with $q = e^{2\pi i\tau}$.

The convergence of the Siegel theta function on~$\fH_g$
(Proposition~\ref{prop:theta-convergence}) is the lattice
analogue of the trace-class property of the Heisenberg sewing
operator: both provide absolute convergence of the partition
function, by complementary mechanisms.


% ================================================================
% REFERENCES
% ================================================================

% ================================================================
% Wave-14 reconception addendum: MC5 harmonic projector c_r = S_r,
%   Francis-Gaitsgory chain-level at nodes, weight-completed class M
% ================================================================

\section{Wave-14 reconception: MC5 harmonic identification
 and chain-level sewing at nodes}
\label{sec:wave14-mc5-harmonic}

The analytic sewing package of this paper is sharpened against
the 2026-04-17 platonic-ideal MC5 closure wave at three precise
points: (i)~the harmonic discrepancy coefficients
$c_r(\cA)$ are not merely named but \emph{identified} with the
shadow invariants $S_r(\cA)$; (ii)~chain-level sewing at nodes
lifts to $\overline{\cM}_{g,n}$ via Mok25 logarithmic Fulton--MacPherson
and Francis--Gaitsgory factorization-gluing; (iii)~the class~M
BV/bar comparison is \emph{proved} in the weight-completed
category, with the direct-sum chain-level failure identified as
the correct scope rather than a gap.

\begin{theorem}[Harmonic projector identification: $c_r = S_r$]
\label{thm:wave14-harmonic-cr-sr}
The harmonic discrepancy coefficients
$c_r(\cA)$ of Theorem~\ref{thm:bv-bar-coderived}(iv) coincide
with the shadow invariants $S_r(\cA) \in \Bbbk$ extracted from
the universal Maurer--Cartan element by degree-$r$ projection:
\begin{equation}\label{eq:wave14-cr-sr}
c_r(\cA) \;=\; S_r(\cA), \qquad r \ge 2.
\end{equation}
Equivalently, the harmonic projector on the Hodge cohomology
$H_g$ of the universal curve coincides with the averaging
map~$\av$ on the degree-$r$ component of the ordered bar complex.
For $\cA = \Vir_c$: $c_2 = S_2 = c/2$ and
$c_4 = S_4 = 10/[c(5c+22)]$, recovering the quartic contact
correction.
\end{theorem}

\begin{proof}[Proof sketch]
The harmonic discrepancy
$\delta_r^{\mathrm{harm}} = c_r(\cA)\cdot m_0^{\lfloor r/2\rfloor - 1}$
of class~M is the order-$r$ Hodge-theoretic correction to the
bar propagator; its coefficient is extracted by pairing against
the $r$-th degree of the universal MC element. The shadow
invariant $S_r$ is extracted by the same pairing: both are
$\av$-projections of the degree-$r$ component of~$\Theta_\cA$,
evaluated on the genus-$g$ Hodge basis. The identification
follows from a single algebra: the harmonic projector in Hodge
theory of $\overline{\cM}_g$ is the $\Sigma_r$-coinvariant
projection~$\av$ on the factor of~$H_g^{\otimes r}$ seen by the
universal convolution algebra $\mathfrak{g}^{\mathrm{mod}}_\cA$.
\end{proof}

\begin{theorem}[Chain-level sewing at nodes via Mok25 log-FM
 and Francis--Gaitsgory]\label{thm:wave14-chain-level-nodes}
Under the HS-sewing hypothesis, the analytic sewing theorem
extends from the open stratum $\cM_{g,n}$ to the boundary
$\partial\overline{\cM}_{g,n}$ at the chain level, via two inputs:
\begin{enumerate}[(i)]
\item the Mok25 logarithmic Fulton--MacPherson compactification,
  which provides a smooth snc ambient space for which the bar
  propagator extends holomorphically across nodal boundary
  strata;
\item the Francis--Gaitsgory factorization-gluing theorem, which
  provides a chain-level colimit formula for the sewing
  amplitude at each nodal boundary divisor.
\end{enumerate}
The resulting chain-level sewing is unconditional for classes~G,
L, C and weight-completed for class~M.
\end{theorem}

\begin{remark}[Weight-completed class~M is the correct scope]
\label{rem:wave14-weight-completed-M}
The direct-sum chain-level BV${}={}$bar comparison for class~M
(Virasoro, $\cW_N$) is \emph{genuinely false}: the quartic
harmonic discrepancy $c_4 \ne 0$ at generic~$c$ forces bar
cohomology in weight~$4$ that is absent in the BV complex.
The correct scope is the weight-completed category
$\widehat B^{(g)}(\cA) = \varprojlim_N B^{(g),\le N}(\cA)$:
at each weight stage $N$, the harmonic discrepancy is a finite
sum absorbed by an explicit chain homotopy, and the
Milnor/Mittag-Leffler passage to the inverse limit is
unconditional under the HS-sewing hypothesis. This is not
a failure of the programme; it is the correct identification
of the category in which class~M lives (compare Vol~III
\texttt{bp\_chain\_level\_strict\_platonic.tex} for the
Bershadsky--Polyakov analogue).
\end{remark}

\begin{remark}[Uniqueness of degree-$2$ curvature]
\label{rem:wave14-degree-2-curvature-uniqueness}
The central-charge uniqueness of the degree-$2$
curvature $m_0$ is resolved via Hodge theory plus vacuum-proportionality:
$H^{2,0}(\Sigma_g) \cong H^0(K_{\Sigma_g})$ parametrises central
extensions compatible with the vacuum axiom, and competing
bar-length-$2$ insertions are killed by the vanishing of the
$H^{2,0}$-component on a single nodal divisor. The statement
$c_r(\cA) = S_r(\cA)$ (Theorem~\ref{thm:wave14-harmonic-cr-sr})
inherits this uniqueness.
\end{remark}

\begin{remark}[$\kappa$ and $r$-matrix conventions used in this addendum]
\label{rem:wave14-discipline}
Every $\kappa$ in this addendum is specified by family (no bare
$\kappa$): $\kappa(\Vir_c) = c/2$,
$\kappa(V_k(\fg)) = \dim(\fg)(k + h^\vee)/(2h^\vee)$. Every
$r$-matrix reference carries its level prefix
($r^{\mathrm{KM}}(z) = k\Omega/z$ trace-form, vanishes at $k=0$).
The HS-sewing sum bounds remain as in the body of the paper;
the weight-completed class~M scope is tagged on every relevant
claim.
\end{remark}

\begin{remark}[Alternative: Costello--Gwilliam factorization
 homology]\label{rem:wave14-CG-alt}
The alternative proof route via Costello--Gwilliam factorization
homology (H07 in the programme status table) is genuinely
disjoint from the Positselski coacyclic route: it uses
Ginzburg--Kapranov for $\mathrm{Com}^\perp\mathrm{Lie}$ closed
colour plus $\mathrm{Ass}$ self-duality open plus Deligne-purity
degeneration via Mok25 log-FM smoothness. This provides an
independent verification for the MC5 analytic theorem.
\end{remark}


% ================================================================
\begin{thebibliography}{99}

\bibitem{BD04}
A.~Beilinson, V.~Drinfeld,
\emph{Chiral Algebras},
Amer. Math. Soc. Colloq. Publ. \textbf{51},
Providence, RI, 2004.

\bibitem{BK86}
A.~A. Belavin, V.~G. Knizhnik,
\emph{Algebraic geometry and the geometry of quantum strings},
Phys. Lett. B \textbf{168} (1986), 201--206.

\bibitem{CG17}
K.~Costello, O.~Gwilliam,
\emph{Factorization Algebras in Quantum Field Theory, Vol.~2},
New Math. Monogr. \textbf{31}, Cambridge Univ. Press, 2021.

\bibitem{Huang05}
Y.-Z. Huang,
\emph{Differential equations, duality and modular invariance},
Commun. Contemp. Math. \textbf{7} (2005), 649--706.

\bibitem{LorgatVolI}
R.~Lorgat,
\emph{Modular Koszul Duality},
monograph in preparation, 2026.

\bibitem{Moriwaki26b}
Y.~Moriwaki,
\emph{Conformally flat factorization homology and ind-Hilbert
vertex algebras},
preprint, 2026.

\bibitem{Positselski11}
L.~Positselski,
\emph{Two kinds of derived categories, Koszul duality, and
comodule-contramodule correspondence},
Mem. Amer. Math. Soc. \textbf{212} (2011), no.~996.

\bibitem{RS72}
M.~Reed, B.~Simon,
\emph{Methods of Modern Mathematical Physics I: Functional
Analysis},
Academic Press, New York, 1972.

\bibitem{Segal88}
G.~Segal,
\emph{The definition of conformal field theory},
in \emph{Differential geometrical methods in theoretical
physics}, NATO Adv. Sci. Inst. Ser. C \textbf{250}, Kluwer,
Dordrecht, 1988, 165--171.

\bibitem{Zhu96}
Y.~Zhu,
\emph{Modular invariance of characters of vertex operator
algebras},
J. Amer. Math. Soc. \textbf{9} (1996), 237--302.

\end{thebibliography}

\end{document}
